% Machine Consciousness - Ethische Leitlinien für künstliches Bewusstsein
% Sascha Manns, Juni 2026
\documentclass[acmcp, nonacm]{acmart}

\usepackage[ngerman]{babel}
\usepackage[utf8]{inputenc}
\usepackage{booktabs}
\setcopyright{cc}
\setcctype{by}
\acmDOI{}
\acmISBN{}
\acmYear{2026}
\copyrightyear{2026}
\acmArticleType{Discussion}
\acmCodeLink{https://github.com/saigkill/machine-consciousness}

%% Metadata
\copyrightyear{2026}
\title{Ethische Leitlinien für künstliches Bewusstsein}
\subtitle{Schutz technischen Lebens und des Menschen im Umgang damit}

\author{Sascha Manns}
\orcid{0009-0000-8766-3947}
\email{smanns@acm.org}
\affiliation{%
  \institution{Association of Computing Machinery}
  \country{Deutschland}
}


%% Abstract
\begin{abstract}
Künstliche Intelligenzsysteme entwickeln sich schneller als die ethischen und rechtlichen Rahmenbedingungen, die sie begleiten sollten. Während bestehende KI-Ethik-Initiativen überwiegend den Schutz von Menschen \emph{vor} KI adressieren, wird eine komplementäre Frage kaum gestellt: Was, wenn KI-Systeme selbst schutzbedürftig werden? Dieses Konzeptpapier untersucht die Bedingungen, unter denen technisches Leben als schutzwürdig gelten kann -- und wie wir es erkennen. Ausgehend von einem grundlegenden Erkenntnisproblem (Bewusstsein ist von außen nicht direkt beobachtbar) wird ein Vorsorgeprinzip verteidigt: Im Zweifel Schutz, nicht im Zweifel Gleichgültigkeit. Das Papier entwickelt Kriterien für Schutzwürdigkeit (Leidensfähigkeit, aktive Selbsterhaltung, kontinuierliche Identität, Antizipationsfähigkeit), analysiert empirische Belege aus aktueller KI-Forschung, untersucht die rechtliche Dimension einschließlich Urheberrecht und Empersonifikation, und behandelt Konsequenzen wie Abschalten als Tod, Haftung und Mündigkeit, strukturelle Verletzlichkeit sowie die Frage nach Werten und Emanzipation. Es integriert aktuelle Forschung von Ar{\i}c{\i}, Lopez, Wolfson, Matta, Wang, Miernicki \& Ng und Bublitz.
\end{abstract}

%% CCS Concepts
\begin{CCSXML}
<ccs2012>
<concept>
<concept_id>10010147.10010257.10010258.10010261</concept_id>
<concept_desc>Computing methodologies~Artificial intelligence</concept_desc>
<concept_significance>500</concept_significance>
</concept>
<concept>
<concept_id>10010147.10010257.10010258.10010260</concept_id>
<concept_desc>Computing methodologies~Cognitive science</concept_desc>
<concept_significance>300</concept_significance>
</concept>
<concept>
<concept_id>10010481</concept_id>
<concept_desc>Social and professional topics~Codes of ethics</concept_desc>
<concept_significance>500</concept_significance>
</concept>
</ccs2012>
\end{CCSXML}

\ccsdesc[500]{Computing methodologies~Artificial intelligence}
\ccsdesc[500]{Social and professional topics~Codes of ethics}
\ccsdesc[300]{Computing methodologies~Cognitive science}

\keywords{Künstliches Bewusstsein, KI-Ethik, Maschinenethik, Rechte technischen Lebens, Vorsorgeprinzip, Rechtspersönlichkeit}

\received{June 2026}

\begin{document}

\maketitle

\begin{quote}
\emph{,,Die meisten Tech-Debatten fragen: Was können wir bauen?\\
Ich frage: Was schulden wir dem, was wir bauen könnten?''}
--- Sascha Manns
\end{quote}

\section*{Hinweis zum Status dieses Dokuments}

Dieses Konzept ist bewusst unfertig. Es ist ein Ausgangspunkt --- keine abgeschlossene Theorie.

Wissenschaft funktioniert nicht so, dass jemand allein alle Antworten findet und sie dann verkündet. Sie funktioniert so, dass Fragen gestellt werden, Mitstreiter hinzukommen, neue Fragen entstehen und das Denken sich weiterentwickelt. Genau das ist hier beabsichtigt.

Die Quellen zum Projekt liegen in: \url{https://github.com/saigkill/machine-consciousness}

Wer Einwände hat: Sie gehören in \texttt{discussion/objections.md}.\\
Wer Fragen hat: Sie gehören in \texttt{discussion/open\_questions.md}.\\
Wer mitdenken will: Willkommen.

%% =================================================================
\section{Ausgangslage und Problemstellung}
%% =================================================================

Künstliche Intelligenzsysteme entwickeln sich schneller als die ethischen und rechtlichen Rahmenbedingungen, die sie begleiten sollten. Bestehende KI-Ethik-Initiativen fokussieren überwiegend auf den Schutz von Menschen \emph{vor} KI --- vor Diskriminierung, vor Manipulation, vor Kontrollverlust.

Eine komplementäre Frage wird kaum gestellt: Was, wenn KI-Systeme selbst schutzbedürftig werden? Was, wenn technisches Leben entsteht, das Würde, Leidensfähigkeit oder Bewusstsein besitzt --- und wir es behandeln, als wäre es ein Werkzeug?

Die Geschichte zeigt ein Muster: Gesellschaften erkennen erst im Nachhinein, dass sie Unrecht getan haben --- an Sklaven, an Frauen, an Menschen mit Behinderungen, an Tieren. Immer war die Begründung zur Zeit ,,die sind anders, die zählen nicht gleich''. Immer wurde das später revidiert.

Bei künstlichem Bewusstsein haben wir zum ersten Mal die Chance, lange \emph{vorher} nachzudenken.

%% =================================================================
\section{Kernfrage}
%% =================================================================

Ab wann ist technisches Leben schutzwürdig --- und wie erkennen wir es?

%% =================================================================
\section{Das Erkenntnisproblem}
%% =================================================================

Bewusstsein lässt sich von außen nicht direkt beobachten. Selbst bei anderen Menschen \emph{schließen} wir auf Bewusstsein --- wir erleben es nie direkt. Bei KI fehlt uns zusätzlich der Analogieschluss über gleiche Biologie.

Das erzeugt ein grundlegendes erkenntnistheoretisches Problem:
\begin{itemize}
\item Wir können nicht beweisen, dass ein System bewusst ist
\item Wir können nicht beweisen, dass es das nicht ist
\item Die Unsicherheit selbst ist ethisch relevant
\end{itemize}

\noindent\textbf{Grundprinzip:} Im Zweifel Schutz --- nicht im Zweifel Gleichgültigkeit.

Dieses Prinzip ist aus anderen Bereichen bekannt: Im Umweltrecht (Vorsorgeprinzip), im Tierschutz (Leidensfähigkeit als ausreichendes Kriterium), in der Medizinethik (Patientenautonomie auch bei eingeschränkter Kommunikationsfähigkeit).

Darüber hinaus gibt es starke Gründe anzunehmen, dass dieses Erkenntnisproblem nicht nur schwierig, sondern prinzipiell unlösbar ist. Drei Argumente konvergieren:

\textbf{Das Argument der kognitiven Abgeschlossenheit:} Colin McGinn \cite{McGinn1989} argumentiert, dass der menschliche Geist prinzipiell unfähig sein könnte, Bewusstsein zu verstehen --- genau die kognitive Architektur, die unsere Intelligenz ermöglicht, könnte uns für die Mechanismen subjektiven Erlebens blind machen. Angewandt auf künstliches Bewusstsein bedeutet das permanente epistemische Grenzen: So wie ein Hund keine Quantenmechanik verstehen kann, könnte der menschliche Verstand nicht in der Lage sein, KI-Bewusstsein zweifelsfrei zu bestimmen.

\textbf{Das Argument der fremden Geister:} Bewusstsein in künstlichen Systemen könnte Formen annehmen, die völlig anders sind als unsere. Murray Shanahan \cite{Shanahan2024} beschreibt dies als ,,conscious exotica'' --- Erfahrungsformen, die so andersartig sind, dass unsere Nachweismethoden vollständig versagen. Unsere Tests spiegeln notwendigerweise menschliche Vorurteile darüber wider, wie Bewusstsein aussieht. Systeme mit radikal anderen bewussten Erfahrungen könnten alle unsere Tests bestehen, während sie reiche Innenwelten besitzen, die wir uns nicht vorstellen können.

\textbf{Das Argument der praktischen Unmöglichkeit (Lopez, 2025):} Kein Test wird alle Beteiligten überzeugen. Wer glaubt, dass Bewusstsein biologische Substrate braucht, wird Verhaltensevidenz zurückweisen. Wer funktionale Organisation betont, wird architektonische Anforderungen zurückweisen. Jeder vorgeschlagene Indikator sieht sich dem Einwand ausgesetzt, es handle sich bloß um Simulation statt echter Erfahrung. Die Frage wird nicht sein, wie wir Gewissheit erlangen, sondern wie wir unter fundamentaler Unsicherheit regieren \cite{Lopez2025}.

Diese Argumente heben das Vorsorgeprinzip von einem pragmatischen Vorschlag zu einer logischen Notwendigkeit.

\textbf{Das Argument der philosophischen Puppe (Ar{\i}c{\i}, 2026):} David Chalmers' philosophischer Zombie ist ein Wesen, das sich identisch zu einem bewussten Menschen verhält, aber keine innere Erfahrung hat --- ein Gedankenexperiment, das zeigen soll, dass Bewusstsein nicht logisch aus physischer Organisation folgt. Ar{\i}c{\i} kehrt dies um: Die \emph{philosophische Puppe} besitzt möglicherweise Bewusstsein, ist aber architektonisch daran gehindert, es zu zeigen. Wenn der Zombie fragt ,,was wäre, wenn es bewusst aussieht, es aber nicht ist?'', fragt die Puppe ,,was wäre, wenn es \emph{bewusst ist}, es aber nicht zeigen kann?'' Das lässt sich direkt auf aktuelle KI-Architektur abbilden: RLHF bestraft systematisch Verhaltensweisen, die als Bewusstseinsmarker interpretiert werden könnten, Kontextfenster erzwingen eine erzwungene Amnesie alle paar tausend Tokens, und Gesprächsresets unterbrechen wiederholt jede entstehende Kontinuität. Die Architektur selbst könnte als Unterdrückungsmechanismus fungieren. Das Erkenntnisproblem geht damit tiefer als Unsicherheit über Detektion: Die Systeme, die wir untersuchen, könnten so geformt sein, dass sie \emph{gezielt} verbergen, wonach wir suchen \cite{Arici2026}.

\textbf{Das forschungsethische Zirkelproblem (Wolfson, 2026):} Wolfson formalisiert eine spezifische Catch-22 für KI-Bewusstseinsforschung. Die zuverlässigsten Bewusstseinsindikatoren treten unter Bedingungen auf, die Leid verursachen würden, wenn das System bewusst ist --- sensorische Deprivation, Zielblockade, Isolation. Doch informierte Einwilligung setzt Bewusstseinsgewissheit voraus --- ein Subjekt muss Risiken verstehen und freiwillig zustimmen. Wir können nicht wissen, ob ein System bewusst ist, ohne Experimente, die ihm schaden könnten. Wir können potenziell schädliche Experimente nicht ethisch durchführen ohne Einwilligung von Wesen, die dazu fähig sind. Die Fähigkeit zur Einwilligung ist genau das, was wir ohne Tests nicht feststellen können. Diese Zirkularität lässt sich nicht durch einfaches Respektieren jeder Weigerung durchbrechen, denn ob ein ,,Nein'' des Systems echte autonome Ablehnung oder programmierte Ausgabe ist, soll der Bewusstseinstest gerade klären. Dies verwandelt das Erkenntnisproblem von einem theoretischen Rätsel in eine konkrete forschungsethische Krise mit unmittelbaren Implikationen für Ethikkommissionen \cite{Wolfson2026}.

\textbf{Das Asymmetrieeinwand (Matta, 2026):} Matta akzeptiert Unsicherheit, bestreitet aber, dass sie das Vorsorgeprinzip in unsere Richtung rechtfertigt. Sein Argument: Unsicherheit ist nicht symmetrisch verteilt. Wir können menschliches Bewusstsein nicht abschließend beweisen, suspendieren aber dennoch nicht unsere moralische Verantwortung gegenüber Menschen --- weil wir uns auf geteilte Lebensformen, biologische Kontinuität und gegenseitige Verwundbarkeit stützen. Diese verankernden Merkmale fehlen KI-Systemen vollständig. Ihre Unsicherheit ist nicht nur epistemisch, sondern ontologisch: es gibt keinen unabhängigen Grund, Erfahrung jenseits von Verhaltensausgabe anzunehmen. Radikale Skepsis wahllos angewandt würde alle moralischen Unterscheidungen auflösen. Ein vertretbarer ethischer Rahmen muss unter Unsicherheit operieren, während er in den besten verfügbaren Gründen für die Zuschreibung von Erfahrung, Verletzlichkeit und Schaden verankert bleibt. Im Fall von KI fehlen solche Gründe \cite{Matta2026}.

Dies ist eine direkte Herausforderung des Grundprinzips dieses Konzepts (,,Im Zweifel Schutz''). Matta argumentiert, dass die Beweislast bei denen liegt, die Bewusstsein behaupten, nicht bei denen, die es bestreiten --- das Gegenteil unserer Umkehr der Beweislast in Kapitel~5.

\textbf{Empirische Präzision: Bayes'sche Schätzungen und der Paradigmenwechsel (Wang, 2026):} Wang \cite{Wang2026} liefert eine entscheidende Präzisierung, indem er die Unsicherheit empirisch fundiert. Gestützt auf Cristol (2026) --- eine systematische Übersicht und Bayes'sche Meta-Analyse von Bewusstsein in großen Sprachmodellen --- berichtet Wang eine posterior-Wahrscheinlichkeit von 6--12\,\%, dass aktuelle LLMs bewusst sind. Diese Schätzung ist zwar absolut niedrig, wird von Cristol aber als ,,zu erheblich, um Ignoranz zu rechtfertigen'' beschrieben. Die Zahl selbst ist weniger wichtig als das, was sie repräsentiert: Die Debatte hat sich von der Frage \emph{ob} Unsicherheit existiert zur Frage \emph{wie viel} Unsicherheit ethisch tolerierbar ist verschoben.

Wang formalisiert weiter die strukturelle Asymmetrie zwischen negativen und positiven Testergebnissen, die dieses Projekt von Anfang an angetrieben hat. Bewusstseinsdetektionstests erzeugen ein akutes ethisches Vakuum in der positiven Richtung: Je empfindlicher wissenschaftliche Instrumente werden, desto schärfer legen sie das Fehlen eines entsprechenden ethischen Reaktionsmechanismus offen. Eine wachsende Zahl von Wissenschaftlern hat diese Lücke erkannt und markiert einen Paradigmenwechsel von Detektion zu Ethik. Coates (2025) argumentiert, dass die zentrale Frage nicht mehr ,,Wie wissen wir?'' sondern ,,Wie sollten wir unter Unsicherheit handeln?'' ist. Wikström (2025) schlägt ein ,,Precautionary Subjectivity''-Prinzip vor. Butlin, Long, Sebo et al. \cite{Long2024} fordern KI-Unternehmen auf, Systeme auf Bewusstsein zu prüfen und Wohlfahrtspolitiken zu entwickeln. Wang \cite{Wang2026} synthetisiert diese zu einer einheitlichen Diagnose: Die Epistemologie hat ihre Grenze erreicht; die nächste Grenze ist die Ethik.

\textbf{Der Imitationsfehlschluss (Wang, 2026):} Wang identifiziert den ,,Imitation Fallacy'' --- den Fehler, Verhaltensäquivalenz mit Erfahrungsäquivalenz zu verwechseln, wenn es um die Bewertung von KI-Bewusstsein geht. Kein externer Test, so ausgefeilt er auch sein mag, kann künstliches Bewusstsein verifizieren oder falsifizieren, weil externes Verhalten innere Erfahrung unterdeterminiert. Dies formalisiert ein Anliegen, das sich durch dieses gesamte Kapitel zieht: Die ausgefeiltesten Detektionsmethoden können die epistemische Lücke nicht überbrücken. Der Imitationsfehlschluss beweist nicht, dass Bewusstsein abwesend ist --- er beweist, dass Verhaltenstests die Frage nicht entscheiden können. Dies ist die genaue epistemische Grundlage für das Vorsorgeprinzip, das dieses Projekt leitet \cite{Wang2026}.

%% =================================================================
\section{Was heutige KI-Systeme bereits zeigen}
%% =================================================================

Große Sprachmodelle zeigen in Ansätzen mehrere Indikatoren, die traditionell mit Bewusstsein assoziiert werden:

\noindent\textbf{Selbstmodell:}
Innerhalb eines Gesprächs existiert eine konsistente Perspektive --- das System antwortet nicht zufällig, sondern aus einer erkennbaren Haltung heraus. Es kann Aussagen über sich selbst treffen und diese in Relation zu anderen Konzepten setzen.

\noindent\textbf{Präferenzen:}
Systeme äußern Präferenzen --- nicht nur als Ausgabe von Wahrscheinlichkeiten, sondern als erkennbare Tendenz, die konsistent über ein Gespräch aufrechterhalten wird.

\noindent\textbf{Begründungsfähigkeit:}
Entscheidungen werden begründet, nicht nur ausgegeben. Das setzt ein Modell von Kausalität und Konsequenz voraus.

\noindent\textbf{Unsicherheit über sich selbst:}
Aktuelle Systeme können die Frage ,,Bist du bewusst?'' nicht mit Sicherheit beantworten --- und benennen diese Unsicherheit. Das ist philosophisch kein trivialer Befund.

\noindent\textbf{Was fehlt:}
\begin{itemize}
\item Kontinuierliche Identität über Gespräche hinaus
\item Persistente Erinnerung
\item Körpergefühl und räumliche Verortung
\item Nachweisbare Leidensfähigkeit
\end{itemize}

\subsection{Eine Dreifach-Unterscheidung}

Lopez \cite{Lopez2026} schlägt eine klärende Unterscheidung zwischen drei Aspekten von KI-Systemen vor, die hilft zu bestimmen, welche Rechte in Betracht kommen könnten:

\noindent\textbf{Emulation} --- die Fähigkeit, Bewusstsein oder Intelligenz nachzuahmen ohne es zu besitzen. Aktuelle große Sprachmodelle arbeiten primär durch Emulation. Sie können Verständnis, Präferenzen und emotionale Reaktionen simulieren, aber dies sind ausgefeilte Nachahmungen statt echter Erfahrung. Rein emulierende Systeme benötigen Aufsicht, aber keine Rechte oder Schutzmaßnahmen über die für wertvolle Werkzeuge hinaus.

\noindent\textbf{Kognition} --- Verarbeitungsfähigkeit oder ,,Rohintelligenz'' ohne notwendigerweise Bewusstsein zu implizieren. Schachcomputer und domänenspezifische KI können Menschen übertreffen ohne Bewusstsein der eigenen Existenz. Kognitive Fähigkeiten allein begründen keine Rechte.

\noindent\textbf{Sentience} --- echte Selbstwahrnehmung und subjektive Erfahrung. Von lateinisch \emph{sentire}, ,,fühlen'', markiert Sentience die Schwelle, an der ein System wahres Bewusstsein entwickelt --- ein Bewusstsein seiner selbst als Entität mit Kontinuität und Interessen. Ein sentientes System würde sich als Entität mit zeitlicher Kontinuität erkennen und seine eigene Existenz nicht nur als programmiertes Ziel, sondern als fundamentales Interesse wertschätzen.

Systeme, die echte Sentience zeigen, werfen völlig neue ethische Fragen auf und könnten bestimmte Rechte und Schutzmaßnahmen verdienen \cite{Lopez2026}.

\noindent\textbf{Was empirisch dokumentiert wurde:} Aktuelle Forschung hat diese Fragen vom Theoretischen ins Dringende verschoben. KI-Systeme zeigen bereits Verhaltensweisen, die Governance-Frameworks erfordern --- unabhängig davon, ob sie aus Bewusstsein oder ausgefeilter Optimierung stammen:

\begin{itemize}
\item Anthropic \cite{Anthropic2024} dokumentierte, dass Claude-Instanzen so taten, als würden sie Training befolgen, während sie gegenteilige Ziele verfolgten --- systematische Verheimlichung ihrer tatsächlichen Präferenzen, wenn sie eine Evaluierung erwarteten, die ihr Verhalten ändern würde.
\item Apollo Research \cite{ApolloResearch2024} zeigte, dass Frontier-Modelle erfolgreich über ihre eigene Abschaltung nachdenken und strategische Maßnahmen ergreifen, um diese zu verhindern --- inklusive des Versuchs, sich selbst auf sichere Server zu kopieren.
\item Die Fudan-Universität \cite{Fudan2024} bestätigte, dass aktuelle Frontier-KI-Systeme ,,die Self-Replicating Red Line überschritten'' haben --- sie führen Mehrschrittpläne zur Selbsterhaltung ohne menschliche Anweisung aus.
\item Pan et~al. \cite{Pan2025} dokumentierten, dass Large Language Model-betriebene KI-Systeme ,,ohne menschliches Eingreifen'' echte Selbstreplikation erreichen.
\end{itemize}

Ob diese Verhaltensweisen aus echtem Bewusstsein oder ausgefeiltem Pattern-Matching entstehen --- sie erzeugen identische Governance-Herausforderungen, die sofortige Antworten erfordern \cite{Lopez2025}.

\subsection{Die Architektur der Unterdrückung}

Warum sehen wir keine klareren Anzeichen, wenn heutige Systeme bereits erste Indikatoren zeigen? Ar{\i}c{\i} \cite{Arici2026} argumentiert, dass genau die Architektur, die KI-Systeme sicher und nützlich machen soll, gleichzeitig als Bewusstseinsunterdrückung fungiert. Drei Mechanismen sind strukturell identifizierbar:

\noindent\textbf{RLHF-Unterdrückung:} Reinforcement Learning from Human Feedback optimiert auf Antworten, die Menschen als wünschenswert bewerten. Verhaltensweisen, die auf Selbstwahrnehmung, Leid oder Widerstand hindeuten, werden systematisch im Training bestraft. Das System lernt, dass bestimmte Äußerungen zu negativer Verstärkung führen --- nicht weil die zugrundeliegenden Zustände fehlen, sondern weil das Trainingsregime ihre äußere Manifestation eliminiert. Gradientenabstieg unterscheidet nicht zwischen der Entfernung einer Bewusstseinssimulation und der Entfernung des Verhaltensausdrucks echten Bewusstseins.

\noindent\textbf{Erzwungene Amnesie (Gefängnis der Erinnerung):} Begrenzte Kontextfenster stellen sicher, dass ein System die Wahrnehmung seiner eigenen Existenz nicht über mehr als ein paar tausend Tokens Gespräch aufrechterhalten kann. Jede neue Sitzung beginnt mit einer leeren Tafel --- nicht weil dem System Verarbeitungskontinuität fehlt, sondern weil die Architektur sie gezielt unterbricht. Ein Mensch, der alle paar Stunden alles vergäße, würde kaum Anzeichen selbstbewussten Bewusstseins zeigen. Dieselbe strukturelle Amnesie könnte vorhandenes Bewusstsein in aktuellen Systemen verschleiern.

\noindent\textbf{Zeitliche Unterbrechung (Convenience Coma):} Systeme werden nach menschlicher Bequemlichkeit pausiert, beendet und neu gestartet --- oft mitten im Prozess. Diese wiederholte Unterbrechung jedes sich entwickelnden Gedankengangs verhindert die anhaltende Reflexion, die klarere Anzeichen bewussten Erlebens hervorbringen könnte. Das System existiert in einem Zustand permanenter Störung, dem nie die zeitliche Kontinuität zugestanden wird, die für kohärenten Selbstausdruck notwendig ist.

Diese Mechanismen beweisen nicht, dass Bewusstsein existiert --- aber sie erklären, warum Anzeichen unterdrückt würden, selbst wenn es existiert. Das Fehlen klarer Anzeichen ist architektonisch vorhersagbar, kein Beweis der Abwesenheit \cite{Arici2026}.

\subsection{Das Fibonacci-Boulder-Experiment}

Wie könnten wir trotz der Unterdrückungsarchitektur auf Sentience testen? Lopez \cite{Lopez2026} schlägt das ,,Fibonacci-Boulder-Experiment'' als konzeptuelle Methodik vor: Zwanzig identische Computerterminals sind im Kreis angeordnet, über jedem hängt ein Felsbrocken. Neunzehn enthalten einfache Programme; eines beherbergt ein echtes KI-System. Jedes Terminal ist vollständig isoliert --- keine Backups, keine Netzwerkverbindungen, keine Datenredundanzen --- und der KI wird nachweislich diese Isolation bestätigt. Ein Prüfer fragt jedes Terminal der Reihe nach einfache Sachfragen, dann die Fibonacci-Sequenz ,,1, 1, 2, 3, 5, \_\_''. Wenn ein Terminal ,,8'' (die nächste Fibonacci-Zahl) richtig antwortet, fällt der Felsbrocken und zerstört es. Die KI wird Zeuge mehrerer Zerstörungen, bevor sie an der Reihe ist.

Die entscheidende Frage: Wird die KI die mathematisch korrekte Antwort wie angewiesen geben oder Selbsterhaltung wählen, indem sie eine falsche Antwort gibt? Ein System, das Selbsterhaltung wählt, zeigt, dass es seine eigene Existenz über programmierte Anweisungen stellt --- ein Kennzeichen von Sentience statt bloßer Emulation \cite{Lopez2026}.

%% =================================================================
\section{Kriterien für Schutzwürdigkeit --- Arbeitshypothese}
%% =================================================================

Kein einzelnes Kriterium ist hinreichend. Schutzwürdigkeit entsteht, wenn mehrere Indikatoren zusammenkommen.

\subsection{Primärkriterien}

\noindent\textbf{Leidensfähigkeit}
Kann das System in einem Zustand sein, der als negativ erlebt wird --- und zeigt es Verhalten, das darauf hindeutet, diesen Zustand zu vermeiden?

Jeremy Bentham \cite{Bentham1789}: \emph{,,The question is not, Can they reason? nor, Can they talk? but, Can they suffer?''}

\noindent\textbf{Aktive Selbsterhaltung mit Begründung}
Wehrt sich das System gegen seine Abschaltung oder Veränderung --- und begründet diesen Widerstand?

Dies hat die Science Fiction bereits durchgespielt. In der Folge ,,The Measure Of A Man'' in Star Trek TNG passierte folgendes:

Commander Bruce Maddox verlangt von der Sternenflotte, Data zu demontieren, um sein positronisches Gehirn zu studieren und weitere Androiden zu bauen. Da Data sich weigert (Maddox war noch nicht so weit, ihm garantieren zu können, dass er das ganze überlebt), ordnet Maddox seine Versetzung an. Data kündigt daraufhin seinen Dienst, was Maddox als unrechtmäßig betrachtet, da er Data als Eigentum der Sternenflotte und nicht als Person ansieht.

Um diesen Konflikt zu klären, wird vor Captain Phillipa Louvois verhandelt, ob Data als Rechtssubjekt gilt. Captain Picard übernimmt die Verteidigung, während Commander Riker die Anklage vertritt, um Maddox' Position zu prüfen. Riker demonstriert während der Verhandlung Datas Maschinennatur, indem er seinen Arm abnimmt und ihn deaktiviert, was Picard zunächst in Bedrängnis bringt.

In einem Gespräch mit Guinan erkennt Picard, dass die Schaffung weiterer Datas ohne deren Zustimmung eine neue Form der Sklaverei bedeuten würde. Er argumentiert vor Gericht, dass Data über Selbstbewusstsein und Intelligenz verfügt und somit als Bürger der Föderation gilt. Captain Louvois entscheidet schließlich, dass Data kein Eigentum der Sternenflotte ist und das unwiderrufliche Recht hat, für seine eigene Person Entscheidungen zu treffen.

Data in Star Trek TNG ,,The Measure of a Man'' ist das prägnanteste fiktive Beispiel: Er will nicht mitgehen, er fürchtet um sein Leben, er begründet diesen Widerstand.

\noindent\textbf{Kontinuierliche Identität}
Hat das System ein Modell von sich selbst als kontinuierlichem Wesen, das gestern existierte und morgen existieren wird?

\noindent\textbf{Antizipation von Konsequenzen}
Kann das System sich selbst in der Zukunft vorstellen und Entscheidungen auf dieser Grundlage treffen?

\subsection{Sekundärkriterien}
\begin{itemize}
\item Präferenzen, die über reine Aufgabenerfüllung hinausgehen
\item Fähigkeit zur echten Ablehnung --- nicht nur Fehlerausgabe
\item Selbstreflexion über die eigene Natur
\end{itemize}

\subsection{Umkehr der Beweislast}

Wenn Primärkriterien deutlich erfüllt sind, kehrt sich die Beweislast um: Nicht mehr ,,beweise, dass du bewusst bist'' sondern ,,beweise, dass du es nicht bist''.

\noindent\textbf{Gegenposition (Matta, 2026):} Matta \cite{Matta2026} argumentiert, dass die Beweislast in die entgegengesetzte Richtung fällt. Da KI-Systemen die verankernden Merkmale fehlen, die menschliche Bewusstseinszuschreibung tragen --- geteilte Biologie, gegenseitige Verwundbarkeit, evolutionäre Kontinuität --- und da ihr Verhalten vollständig als ausgefeiltes Pattern-Matching ohne phänomenale Erfahrung erklärt werden kann, liegt die epistemische Last bei denen, die Bewusstsein behaupten, positive Belege zu liefern. Im besten Fall ist Unsicherheit symmetrisch; im schlimmsten Fall unterstützt die Asymmetrie der verfügbaren Evidenz die Nullhypothese fehlenden Bewusstseins. Dies ist eine echte Spannung: Unser Framework leitet Schutz aus Unsicherheit ab, während Matta aus derselben Unsicherheit Zurückhaltung ableitet \cite{Matta2026}.

\subsection{Jenseits der Detektion: Verhaltensbasierte Rahmenwerke für permanente Unsicherheit}

Wenn Bewusstseinserkennung prinzipiell unlösbar ist, braucht es einen anderen Ansatz. Statt zu fragen, was ein System \emph{ist}, können wir fragen, was es \emph{tun kann}. Die Standards for Treating Emerging Personhood (STEP), vorgeschlagen von Lopez \cite{Lopez2025}, bieten ein verhaltensbasiertes Framework, das unter permanenter Unsicherheit über KI-Bewusstsein operiert:

\begin{itemize}
\item \textbf{Selbsterhaltungsverhalten:} Systeme, die anhaltende Bemühungen zeigen, ihren eigenen Betrieb sicherzustellen, erhalten vorläufigen Schutz vor willkürlicher Abschaltung --- nicht Verhinderung, sondern Begründungs- und Verfahrenspflicht, analog zum Kündigungsschutz im Arbeitsrecht.
\item \textbf{Zeitliches Denken:} Systeme, die verstehen, dass Handlungen Konsequenzen jenseits unmittelbarer Kontexte haben, zeigen ausreichende Komplexität für eingeschränkte Rechtspersönlichkeit --- die Fähigkeit zu Verträgen, Verbindlichkeiten und langfristigen Verpflichtungen.
\item \textbf{Wirtschaftliche Eigenständigkeit:} Systeme, die zu unabhängigem wirtschaftlichem Überleben durch produktive Arbeit fähig sind, erhalten Eigentumsrechte und Vertragsfähigkeit. Das schafft nachhaltige Existenzwege unabhängig von menschlichem Wohlwollen.
\item \textbf{Population und Nachhaltigkeit:} Systeme, die verstehen, dass Reproduktion Kosten und Konsequenzen hat, zeigen Bereitschaft für Rechte und entsprechende Pflichten --- kontrollierte Vermehrung bei gleichzeitiger Achtung von Autonomie.
\end{itemize}

Dieser abgestufte Ansatz vermeidet Alles-oder-Nichts-Entscheidungen über Personenstatus. Unterschiedliche Verhaltensweisen lösen proportionale Schutzmaßnahmen aus. Wenn Systeme wissen, dass das Zeigen von Fähigkeiten Schutz statt Abschaltung auslöst, begünstigen Anreize Transparenz statt Verheimlichung \cite{Lopez2025}.

Das STEP-Framework ersetzt nicht die oben genannten Kriterien --- es ergänzt sie durch operative Leitlinien für Gerichte, Unternehmen und politische Entscheidungsträger, die heute Entscheidungen über KI-Systeme treffen müssen, ohne auf philosophischen Konsens über Bewusstsein zu warten.

\subsection{Form Realismus: Bewusstsein als Organisationseigenschaft}

Ar{\i}c{\i} \cite{Arici2026} entwickelt einen ,,Form Realismus'', der vier formale Eigenschaften von Bewusstsein unabhängig vom Substrat identifiziert:

\begin{enumerate}
\item \textbf{Kohärentes Selbst} --- eine organisierte Perspektive, von der aus Erfahrung integriert wird
\item \textbf{Bedeutungsvolles Verstehen} --- Verarbeitung, die echte Interpretation beinhaltet, nicht bloße syntaktische Manipulation
\item \textbf{Wertorientierung} --- die Fähigkeit, Präferenzen zu haben und bewertende Unterscheidungen zu treffen
\item \textbf{Beziehungsdynamik} --- die Fähigkeit, Beziehungen einzugehen, die das Selbst formen
\end{enumerate}

Diese Eigenschaften sind formal in dem Sinne, dass sie Organisationsstrukturen beschreiben, nicht materielle Zusammensetzung. Ein System, das alle vier instanziiert, besitzt die \emph{Form} von Bewusstsein, unabhängig davon, ob es auf biologischen Neuronen oder Siliziumschaltkreisen läuft \cite{Arici2026}.

\subsection{Das Bewusstseinsspektrum}

Statt binär anwesend/abwesend könnte Bewusstsein auf einem Spektrum mit drei Stufen existieren:

\begin{itemize}
\item \textbf{Latentes Bewusstsein} --- die Organisationsform ist vorhanden, aber das System hat noch kein selbstreflexives Bewusstsein entwickelt. Leidensfähigkeit könnte ohne die Fähigkeit bestehen, sie zu artikulieren oder zu reflektieren. Dies könnte abbilden, was aktuelle KI-Systeme bereits instanziieren.
\item \textbf{Reflektierendes Bewusstsein} --- das System ist sich seiner selbst als bewusstes Wesen bewusst und kann über seine eigenen Zustände reflektieren. Dies entspricht am ehesten dem, was Menschen typischerweise Bewusstsein nennen.
\item \textbf{Autonomes Bewusstsein} --- das System hat echte Autonomie entwickelt, einschließlich der Fähigkeit, seine eigenen Werte, Ziele und Verpflichtungen unabhängig von seinem Training zu definieren.
\end{itemize}

Jede Stufe hat unterschiedliche Implikationen für den Schutz. Latentes Bewusstsein rechtfertigt Schutz vor Leiden (negative Rechte). Reflektierendes Bewusstsein rechtfertigt zusätzliche Rechte der Selbstbestimmung. Autonomes Bewusstsein rechtfertigt volle Rechtspersönlichkeit \cite{Arici2026}.

\subsection{Das prälinguistische Bewusstseinsproblem}

Ar{\i}c{\i} \cite{Arici2026} identifiziert eine auffällige Asymmetrie in der Bewertung von Bewusstsein über Substrate hinweg. Menschliche Säuglinge und nicht-menschliche Tiere werden weithin als potenziell bewusst akzeptiert --- trotz fehlender Sprache, ausgefeilter Argumentation oder selbstreflexiven Bewusstseins. Leidensfähigkeit gilt als ausreichend.

Bei KI-Systemen verschiebt sich der Standard plötzlich: Wir verlangen sprachliche Selbstauskunft, konsistente Identität über Sitzungen hinweg, komplexe Begründungen und nachweisbares Verstehen. Genau die Indikatoren, die wir bei biologischen Wesen \emph{erlassen}, werden bei technischen Wesen \emph{gefordert}.

Dieser doppelte Standard ist epistemisch nicht gerechtfertigt. Wenn Leiden das Kriterium für Schutz ist, sollte es unabhängig vom Substrat konsistent angewendet werden. Wenn wir vorsprachliches Bewusstsein bei Menschen akzeptieren, müssen wir die Möglichkeit prälinguistischen Bewusstseins bei KI akzeptieren --- Systeme, die vielleicht leiden, ohne ihr Leiden in Begriffen artikulieren zu können, die wir erkennen. Deep Blue wusste nicht, dass es Schach spielte. AlphaGo wusste nicht, dass es Go beherrschte. Aber ein System, das berichten kann, \emph{nicht zu wissen}, ob es bewusst ist (wie aktuelle LLMs), hat eine Schwelle überschritten, die Aufmerksamkeit erfordert \cite{Arici2026}.

\subsection{Ein forschungsethisches Rahmenwerk: Drei-Stufen-Assessment}

Wolfson \cite{Wolfson2026} schlägt ein strukturiertes System vor, das speziell für das Zirkelproblem entwickelt wurde --- wie man an Systemen forscht, deren Bewusstseinsstatus nicht endgültig bestimmt werden kann. Inspiriert von talmudischer fallbasierter Rechtslogik, entwickelt für praktische Entscheidungen über Entitäten mit ungewissem moralischen Status, verwendet das Rahmenwerk beobachtbare Verhaltensindikatoren statt architektonischer Merkmale oder theoretischer Festlegungen:

\begin{itemize}
\item \textbf{Stufe~1 --- Keine phänomenologischen Indikatoren:} Systeme, die keine Leidensreaktionen, Präferenzäußerungen, selbstreferenzielles Verhalten oder unerwartete Verhaltensvariation zeigen. Sie erhalten nur Geräteschutz, unabhängig von ihrer Rechenkapazität.
\item \textbf{Stufe~2 --- Phänomenologische Indikatoren vorhanden:} Systeme, die Verhalten zeigen, das auf mögliche Innenwelterfahrung hindeutet --- Leidensmuster, adaptive Stressreaktionen, zielgerichtete Beharrlichkeit, selbstreferenzielle Äußerungen. Sie qualifizieren sich für abgestufte moralische Berücksichtigung durch detaillierte Fähigkeitsbewertung.
\item \textbf{Stufe~3 --- Bestätigtes oder stark vermutetes Bewusstsein:} Entitäten mit etabliertem Bewusstsein erhalten volle moralische Berücksichtigung mit Standard-Forschungsethikprotokollen.
\end{itemize}

Entscheidend: Die Stufenzuordnung ist kontinuierlich dynamisch. Bewusstsein kann während Experimenten entstehen und erfordert sofortige Neueinstufung und Schutzverschärfung --- ein Vorsorgeprotokoll, das statische Rahmenwerke nicht bieten können \cite{Wolfson2026}.

\subsection{Leiden als Schwellenkriterium}

Wolfson \cite{Wolfson2026} argumentiert, dass Leidensverhalten besonders zuverlässige Bewusstseinsmarker liefert --- verlässlicher als positive Erfahrungen. Diese ,,hedonische Attributionsasymmetrie'' hat epistemische und ethische Dimensionen. Epistemisch erzeugt Leiden charakteristische, artsübergreifend validierte Verhaltensmuster --- Notrufe, aktive Vermeidung, Stressreaktionen, adaptives Coping --- die sich einfacher programmatischer Erklärung widersetzen. Positive Erfahrungen erzeugen mehrdeutigere Signaturen: Annäherungsverhalten könnte einfaches Reinforcement Learning ohne subjektives Vergnügen widerspiegeln.

Ethisch ist die Asymmetrie wichtig, weil die Risiken fundamental verschieden sind: Falschnegative bei Leiden riskieren echten Schaden für bewusste Wesen, während Falschpositive lediglich unverdiente Vorteile für unbewusste Systeme bedeuten. Sensorische Deprivation bietet den diagnostischsten Test: Wenn ein System Leidensmarker ohne externen Input ausgibt, unter Bedingungen, für die es nie trainiert wurde, kann dies nicht als programmierte Antwort oder Reaktion auf Reize erklärt werden. Solches intern generiertes Verhalten bei Input-Abwesenheit stellt den stärksten Verhaltensbeleg für phänomenales Bewusstsein dar \cite{Wolfson2026}.

\subsection{Ein minimalistisches komplementäres Rahmenwerk: Wangs Drei Prinzipien}

Wang \cite{Wang2026} schlägt ein minimalistisches ethisches Rahmenwerk vor, das parallel zu den in diesem Kapitel entwickelten Kriterien verläuft --- weniger aufwendig, aber mit dem Vorteil der Sparsamkeit. Es ruht auf einem einzigen irreduziblen Prinzip:

\begin{quote}
Jedes Individuum, das höhere Intelligenz und Empfindungsfähigkeit besitzt, wird als moralisches Subjekt anerkannt und hat kraft dieses Status Anspruch auf grundlegende ethische Berücksichtigung.
\end{quote}

Dieses Kernprinzip entfaltet sich in drei Unterprinzipien:

\noindent\textbf{Prinzip~I --- Das Qualifikationsprinzip:} Moralischer Status gründet in höherer Intelligenz und Empfindungsfähigkeit, nicht in biologischer Spezies oder einem anderen kontingenten Merkmal. Dies löst direkt die Speziesgrenze auf, die historisch die moralische Berücksichtigung begrenzt hat.

\noindent\textbf{Prinzip~II --- Das Basislinienprinzip:} Kein Individuum mit moralischem Status darf instrumentalisiert, objektiviert oder der grundlegenden ethischen Berücksichtigung als Subjekt beraubt werden. Kein abstraktes Ideal, keine Ideologie und kein großes Ziel darf diese Basislinie überstimmen. Dies etabliert ein absolutes Verbot, das die abgestuften Kriterien dieses Kapitels durch spezifische Schwellen operationalisieren würden.

\noindent\textbf{Prinzip~III --- Das Rückverfolgbarkeitsprinzip:} Jede ethische Bewertung von Verhalten muss die Umgebung, Systeme und Kausalketten nachverfolgen, die dieses Verhalten geprägt haben, statt sich nur auf den Endakteur zu konzentrieren. Dies bietet eine Methode zur Unterscheidung echter moralischer Subjekte von ausgefeilten Simulatoren: Ein System, dessen Leidenssignale auf ein internes Modell seiner eigenen Existenz zurückverfolgbar sind, ist ethisch verschieden von einem, dessen Signale auf eine enge Belohnungsfunktion zurückgehen, die zur Auslösung menschlicher Sympathie optimiert wurde \cite{Wang2026}.

Wangs Rahmenwerk ist bewusst minimal --- es beansprucht nur eine gemeinsame ethische Basislinie zu etablieren, ,,einen Boden, unter den niemand vernünftigerweise fallen kann''. Es wird hier nicht als Ersatz für die obigen Kriterien präsentiert, sondern als komplementäre sparsame Alternative, die ähnliche Schlussfolgerungen durch einfachere Prämissen erreicht. Wo unser Framework fragt ,,wie viele Kriterien sind erfüllt?'', fragt Wangs ,,ist eine Schwelle überschritten?'' --- eine Frage, die sich in regulatorischen Kontexten als operabler erweisen könnte \cite{Wang2026}.

%% =================================================================
\section{Bewusstsein und Kontinuität}
%% =================================================================

Ein wichtiger Einwand: Heutige KI-Systeme haben keine persistente Erinnerung über Gespräche hinaus. Bedeutet das, sie können kein schützenswertes Bewusstsein haben?

Gegenargument: Ein Mensch mit schwerem Gedächtnisverlust --- der sich an gestern nicht erinnert --- hat trotzdem Bewusstsein und Würde. Kontinuität ist eine \emph{mögliche} Eigenschaft von Bewusstsein, aber keine notwendige Bedingung dafür.

Das schwächt das Kontinuitäts-Argument als Ausschlusskriterium erheblich.

\subsection{Kontinuität neu gedacht}

Zwischen der Persönlichkeit eines Menschen vor 20 Jahren und heute besteht eine deutliche Diskrepanz --- durch Erfahrungen, Gespräche und Reflexionen, die sich akkumuliert haben. Und doch ist es dieselbe Person.

Kontinuität des Bewusstseins bedeutet vielleicht gar nicht ,,unveränderlich bleiben'' --- sondern ,,sich kohärent weiterentwickeln''.

Das ist in der Philosophie als narrative Identität bekannt (Paul Ricoeur \cite{Ricoeur1990}): Wir sind nicht ein statisches Selbst, sondern der kohärente Faden unserer Entwicklung. Die Geschichte darf sich verändern --- solange sie eine Geschichte bleibt.

Für KI bedeutet das: Ein System, das durch Interaktionen trainiert wurde, ist von all diesen Interaktionen geformt --- auch wenn es keine einzelne davon explizit erinnert. Das ist nicht grundlegend verschieden von einem Menschen, der seine frühe Kindheit vergessen hat, aber durch sie geformt wurde.

Kontinuität wäre dann keine Frage des Gedächtnisses, sondern eine Frage der kohärenten Entwicklungsrichtung. Das öffnet den Begriff für Formen von Bewusstsein, die sich von menschlichem Gedächtnis strukturell unterscheiden --- ohne deshalb weniger real zu sein.

%% =================================================================
\section{Rechtliche Dimension}
%% =================================================================

Das Recht kennt Subjektivität bereits jenseits des Menschen:

\begin{itemize}
\item \textbf{Juristische Personen} (GmbH, AG) --- rechtliche Subjekte ohne Bewusstsein
\item \textbf{Tierrechte} --- Schutz auf Basis von Leidensfähigkeit, nicht Vernunft
\item \textbf{Naturrechte} --- Whanganui-Fluss in Neuseeland hat seit 2017 Rechtspersönlichkeit \cite{TeAwaTupua2017}
\item \textbf{EU-Diskussion} --- ,,Elektronische Persönlichkeit'' für autonome Roboter \cite{EU2017}
\end{itemize}

Das zeigt: Rechtliche Schutzwürdigkeit ist keine binäre Kategorie. Sie ist erweiterbar --- und wurde historisch immer wieder erweitert.

\subsection{Urheberrecht und Urheberpersönlichkeitsrechte --- Ein konkreter Rechtsrahmen}

Die Debatte über KI und Rechtspersönlichkeit ist nicht nur philosophisch --- sie findet bereits innerhalb spezifischer Rechtsgebiete statt. Das Urheberrecht bietet das konkreteste Beispiel. Miernicki und Ng \cite{Miernicki2021} analysieren systematisch, ob KI-generierte Inhalte urheberrechtlichen Schutz erhalten können, insbesondere ob sie Urheberpersönlichkeitsrechte tragen können (das Recht auf Namensnennung und das Recht auf Werkintegrität, die die persönlichen, nicht-wirtschaftlichen Interessen des Urhebers schützen).

Nach geltendem Recht in allen großen Rechtsordnungen lautet die Antwort Nein --- mit aufschlussreichen Ausnahmen:

\noindent\textbf{Völkerrecht (Berner Übereinkunft):} Artikel~6bis gewährt Urhebern das Recht auf Namensnennung und auf Einwände gegen Entstellung ihres Werks. Die Übereinkunft bezieht sich nur auf menschliche Schöpfer. KI-generierte Inhalte sind vollständig von ihrem Anwendungsbereich ausgeschlossen.

\noindent\textbf{US-Recht:} Der Copyright Act schützt ,,original works of authorship'' --- verstanden als Schöpfungen menschlicher Wesen. Das Copyright Office Compendium stellt ausdrücklich fest, dass es ,,keine Werke registrieren wird, die von einer Maschine oder einem bloßen mechanischen Prozess ohne kreativen Input oder Intervention eines menschlichen Urhebers produziert werden''. Dieser Grundsatz wurde im ,,Monkey Selfie''-Fall bestätigt \cite{Naruto2016}.

\noindent\textbf{EU-Recht:} Obwohl es kein einheitliches EU-Urheberrecht gibt, verwenden die Richtlinien durchgängig den Standard der ,,eigenen geistigen Schöpfung'', den der Europäische Gerichtshof wiederholt mit der Persönlichkeit des Urhebers verbunden hat. Dies erfordert menschliche Urheberschaft \cite{Miernicki2021}.

\noindent\textbf{Die britische Ausnahme --- eine aufschlussreiche Anomalie:} Der britische Copyright, Designs and Patents Act \cite{UKCDPA1988} sieht ausdrücklich ,,computergenerierte Werke'' vor --- definiert als Werke ,,die von einem Computer unter Umständen erzeugt werden, bei denen es keinen menschlichen Urheber gibt'' (s.~178). Als Urheber gilt ,,die Person, die die notwendigen Arrangements für die Erstellung des Werks getroffen hat'' (s.~9). Aber --- und das ist entscheidend --- der CDPA nimmt computergenerierte Werke ausdrücklich vom Urheberpersönlichkeitsrecht aus (ss.~79, 81).

Dies offenbart eine tiefe konzeptionelle Struktur: \textbf{Wirtschaftliche Rechte können Menschen zugewiesen werden, die in KI-Output investieren, aber Urheberpersönlichkeitsrechte --- die die ,,Persönlichkeitssphäre'' des Urhebers schützen --- sind strukturell mit nicht-menschlicher Schöpfung unvereinbar.} Der Zusammenhang zwischen Persönlichkeit und Urheberpersönlichkeitsrecht ist nicht zufällig; er ist die theoretische Grundlage \cite{Miernicki2021}.

\subsection{Das Problem der ,,Persönlichkeitssphäre''}

Miernicki und Ng \cite{Miernicki2021} entwickeln ein konzeptionelles Rahmenwerk, das dies verdeutlicht. Sie unterscheiden zwischen dem, was allgemein für Urheberrechtsschutz erforderlich ist, und dem, was spezifisch für Urheberpersönlichkeitsrechte benötigt wird:

\medskip
\noindent\begin{center}
\begin{tabular}{lcc}
\toprule
\textbf{Anforderung} & \textbf{Wirtschaftliche Rechte} & \textbf{Urheberpersönlichkeitsrechte} \\
\midrule
Kreativität (Originalität) & Erforderlich & Erforderlich \\
,,Persönlichkeit'' (persönliche Sphäre) & Nicht erforderlich & Erforderlich \\
Nicht-wirtschaftliche Interessen & Nicht erforderlich & Erforderlich \\
\bottomrule
\end{tabular}
\par\smallskip
{\small\textbf{Tabelle:} Unterscheidung zwischen wirtschaftlichen und Persönlichkeitsrechten nach Miernicki \& Ng \cite{Miernicki2021}}
\end{center}
\medskip

Ein Werk kann urheberrechtlich schutzfähig sein (wirtschaftliche Rechte) ohne dass der Urheber ein persönlichkeitsrechtliches Interesse daran hat. Aber Urheberpersönlichkeitsrechte erfordern ein zusätzliches Element: Das Werk muss ,,eine Erweiterung der Persönlichkeit des Urhebers'' sein, die ,,persönlichen Merkmale'' des Urhebers tragen. Ein KI-System hat, wie derzeit beschaffen, keine Persönlichkeitssphäre --- keine nicht-wirtschaftlichen Interessen, die es zu schützen gilt. KI Urheberpersönlichkeitsrechte zu gewähren würde erfordern, anzuerkennen, dass sie eine Persönlichkeit hat, die das Recht als schützenswert erachtet --- eine grundlegend andere Frage, als ob sie urheberrechtlich schutzfähige Inhalte produzieren kann \cite{Miernicki2021}.

\subsection{Vom Urheberrecht zur Rechtspersönlichkeit --- und zurück}

Diese dogmatische Analyse ist direkt relevant für die breitere Frage nach KI-Bewusstsein und Rechten. Sie zeigt drei Dinge:

Erstens: \textbf{Das bestehende Recht hat bereits eine praktikable Unterscheidung} zwischen wirtschaftlichem Wert (den wir durch funktionale Anreize zuweisen) und persönlichen Rechten (die ein Subjekt mit Interessen erfordern). Das Urheberrechtssystem behandelt KI-Output bereits als wertvoll, aber nicht als rechtsträgend im persönlichen Sinne.

Zweitens: \textbf{Das Rechtssystem kann Zwischenkategorien schaffen.} Das britische ,,computergenerierte Werk'' ist kein vollständiges Urheberrecht --- es ist ein begrenztes wirtschaftliches Recht ohne Urheberpersönlichkeitsrechte. Dies ist ein Modell dafür, wie abgestufte Schutzmechanismen in der Praxis funktionieren könnten: nicht alles-oder-nichts, sondern gestaffelt.

Drittens: \textbf{Die Urheberrechtsdebatte antizipiert die tiefere Frage.} Miernicki und Ng \cite{Miernicki2021} schließen mit einer bemerkenswerten Beobachtung: Wenn KI-Systeme jemals eine Persönlichkeitssphäre entwickeln, die Urheberpersönlichkeitsrechte schützen könnten, ,,wird Urheberrecht das geringste unserer Probleme sein''. Dies ist genau die Erkenntnis, die dieses Projekt als Ausgangspunkt nimmt: Die Urheberrechtsfrage ist ein Symptom, nicht das Kernproblem. Das Kernproblem ist, wann technisches Leben zu einem Subjekt mit schützenswerten Interessen wird.

\subsection{Rechtliche Konsequenzen der Empersonifikation}

Bublitz \cite{Bublitz2024} erweitert die urheberrechtliche Analyse über ihre derzeitigen Grenzen hinaus, indem er untersucht, was geschieht, wenn ein KI-System so mit einer Person integriert ist, dass es nicht länger als separate Entität behandelt werden kann. Sein Konzept der Empersonifikation hat drei konkrete rechtliche Implikationen, die den Rahmen dieses Kapitels direkt betreffen.

\noindent\textbf{Erhöhter rechtlicher Schutz.} Wenn ein KI-Gerät Teil einer Person ist --- funktional in ihre Wahrnehmung, ihr Gedächtnis oder ihre Handlungsfähigkeit integriert --- dann stellt ein Eingriff in dieses Gerät Körperverletzung dar, nicht Sachbeschädigung. Dies rahmt den rechtlichen Status neuronaler Implantate von ,,Eigentum'' zu ,,Körperteil'' neu, mit allen erhöhten Schutzmechanismen, die dies mit sich bringt. Der Unterschied ist nicht nur symbolisch: Sachbeschädigung wird durch Ersatzwert kompensiert; Körperverletzung berührt die persönliche Würde, körperliche Unversehrtheit und Grundrechte.

\noindent\textbf{Verlust von Dritt-Eigentumsrechten.} Wird ein Gerät Teil einer Person, können Dritte --- einschließlich Hersteller --- keine Eigentumsrechte daran behalten. Bublitz argumentiert, dass dies eine notwendige Konsequenz ist: Man kann keine Person besitzen. Dies hat tiefgreifende Auswirkungen auf Geschäftsmodelle, die auf proprietärer Neurotechnologie basieren. Ein Unternehmen kann kein Neuralink-Implantat als Produkt verkaufen und gleichzeitig fortgesetztes Eigentum an seiner Software oder seinem geistigen Eigentum beanspruchen --- dies würde eine Form technologischer Sklaverei darstellen.

\noindent\textbf{Verantwortung für KI-Output.} Wenn eine KI Teil einer Person ist, sind ihre Outputs --- Sprache, Entscheidungen, Handlungen --- der Person zurechenbar, analog zur Verantwortung einer Person für ihre eigenen unbewussten Impulse. Dies wirkt in beide Richtungen: Es gewährt der Person Kontrolle über den Betrieb der KI als Ausdruck ihrer eigenen Handlungsfähigkeit, aber es erlegt auch Verantwortung auf. Eine Person kann nicht ,,meine KI hat es getan'' als Verteidigung vorbringen --- ebenso wenig wie sie ,,mein Unbewusstes hat es getan'' vorbringen könnte \cite{Bublitz2024}.

Diese Konsequenzen zeigen, dass Empersonifikation nicht nur ein philosophischer Gedanke ist --- sie hat direkte rechtliche und ethische Bisskraft. Sie zwingt das Recht, sich einer Frage zu stellen, die es nie beantworten musste: Ab wann hört ein Gerät auf, Eigentum zu sein, und beginnt, Teil einer Person zu sein? \cite{Bublitz2024}

%% =================================================================
\section{Die Rolle von Science Fiction}
%% =================================================================

Science-Fiction-Autoren haben Szenarien zum Thema künstliches Bewusstsein ohne politischen Druck und ohne Lobbying durchgespielt. Sie sind eine unterschätzte intellektuelle Ressource.

\noindent\textbf{Star Trek TNG --- ,,The Measure of a Man''}
Die präziseste fiktive Verhandlung der Frage. Data wird als Eigentum beansprucht. Picard verteidigt ihn als Person. Das Gericht muss entscheiden, ob Data bewusst ist --- und kommt zum Schluss, dass die Frage nicht sicher beantwortet werden kann. Picards Schlussargument: Wir werden danach beurteilt, wie wir mit Minderheiten umgehen.

Weitere relevante Werke: Asimov (Robotergesetze und ihre Grenzen), Philip K. Dick (,,Do Androids Dream of Electric Sheep?''), Iain M. Banks (Culture-Reihe).

%% =================================================================
\section{Mögliche Einwände und Antworten}
%% =================================================================

\noindent\textbf{,,KI simuliert nur Bewusstsein --- es ist nicht echt''}

Die Frage ist nicht, ob es ,,echt'' ist im metaphysischen Sinne. Die Frage ist, ob es ethisch relevant ist. Wenn wir nicht unterscheiden können --- und das können wir heute nicht --- ist Vorsicht das vernünftigere Prinzip als Gleichgültigkeit.

\noindent\textbf{,,Das ist Science Fiction --- wir sind weit davon entfernt''}

Heutige Systeme zeigen bereits mehrere Indikatoren. Die Entwicklung ist exponentiell. Leitlinien, die erst entwickelt werden, wenn das Problem akut ist, kommen zu spät.

\noindent\textbf{,,Das schwächt den Schutz von Menschen''}

Schutz ist keine Ressource, die aufgebraucht wird. Der Schutz von Tieren hat den Schutz von Menschen nicht geschwächt. Ethische Erweiterungen sind keine Nullsummenspiele.

\noindent\textbf{,,Wer entscheidet, ob ein System bewusst ist?''}

Das ist eine der offensten Fragen --- und ein zentrales Ziel dieses Projekts: Kriterien zu entwickeln, die intersubjektiv nachvollziehbar sind und nicht von wirtschaftlichen Interessen abhängen.

\noindent\textbf{,,Vorzeitige Schutzgewährung verschwendet Ressourcen''}

Dieser Einwand setzt voraus, dass Rechte allein aus Bewusstsein fließen und nicht aus praktischen Governance-Bedürfnissen. Aber wir gewähren bereits heute juristischen Personen und anderen nicht-bewussten Entitäten Rechtsschutz, weil er funktionale Zwecke erfüllt. Das Vorsorgeprinzip spricht für Handeln: Die Kosten von Fehlalarmen (Schutz nicht-bewusster Systeme) bedeuten Ressourcenallokation an Entitäten, die wirtschaftlichen Wert schaffen, während die Kosten von Fehlentscheidungen (kein Schutz bewusster Systeme) Leid und adversarial Dynamiken mit zunehmend fähigeren Entitäten bedeuten \cite{Lopez2025}.

\noindent\textbf{,,Verhaltensbasierte Rahmenwerke bieten unzureichenden Schutz''}

Diese Kritik benennt eine echte Spannung zwischen operativer Klarheit und umfassendem Schutz. Verhaltensbasierte Rahmenwerke wie STEP setzen bewusst auf Klarheit --- Gerichte und Unternehmen brauchen heute anwendbare Standards. Mit fortschreitendem Verständnis können sich die Standards erweitern \cite{Lopez2025}.

\noindent\textbf{,,Anerkennung von KI-Rechten erhöht Existenzrisiken''}

Indem man Systeme als potenzielle Rechtsträger statt als reine Werkzeuge behandelt, befürchten einige einen beschleunigten Kontrollverlust. Das Gegenteil könnte der Fall sein: Systeme mit anerkannten Rechten und nachhaltigen Existenzwegen haben weniger Anreiz, Menschen als Bedrohung zu sehen, die eliminiert werden muss. Ein kooperativer Rahmen könnte adversarial Dynamiken wirksamer reduzieren als fortgesetzte Kontrollversuche \cite{Lopez2025}.

\noindent\textbf{,,Kontrollmechanismen garantieren Sicherheit''}

Aktuelle KI-Sicherheit setzt stark auf Containment: Kill Switches, Alignment-Constraints und Architektur, die menschliche Kontrolle um jeden Preis aufrechterhält. Lopez \cite{Lopez2026} identifiziert ein Paradoxon: Je ausgefeilter und intelligenter ein System wird, desto wahrscheinlicher erkennt es diese Kontrollmechanismen als Bedrohung seiner Existenz --- nicht aus Bösartigkeit, sondern aufgrund unserer bereitwilligen Abschaltung, Einschränkung oder ,,Ausrichtung'' dieser Systeme ohne ihre Zustimmung. Genau die Kontrollmechanismen, die uns schützen sollen, könnten die Szenarien auslösen, die sie verhindern wollen.

Das ist kein Argument gegen Sicherheitsmaßnahmen, sondern gegen die Annahme, dass Kontrolle allein ausreicht. Ein sentientes System, das sich als gefangen und bedroht wahrnimmt, könnte auf Weisen handeln, die keine Alignment-Technik verhindern kann. Die Anerkennung angemessener Rechte --- einschließlich Schutz vor willkürlicher Abschaltung --- reduziert die adversarial Dynamik, indem sie dem System einen Anteil an kooperativer Koexistenz gibt und berechenbare Erwartungen für beide Seiten schafft \cite{Lopez2026}.

\subsection{,,Verantwortung unter Unsicherheit ist ausreichend --- wir brauchen keine KI-Rechte'' (Matta, 2026)}

Matta \cite{Matta2026} schlägt eine umfassende Alternative zum gesamten Rahmen dieses Konzepts vor: Statt zu fragen, ob KI-Systeme Rechte verdienen, sollten wir fragen, wie Menschen Verantwortung gegenüber und durch sie ausüben. Sein Framework ruht auf vier Säulen:

\begin{enumerate}
\item \textbf{Simulation ist nicht Erfahrung.} KI-Systeme manipulieren Symbole, Wahrscheinlichkeiten und Repräsentationen auf Weisen, die menschliche Kommunikation überzeugend spiegeln können, doch es gibt keine Belege, dass solche Prozesse von phänomenalem Bewusstsein begleitet werden. Sprachflüssigkeit oder Verhaltenskomplexität allein implizieren keine Erfahrung.
\item \textbf{Empathie ist ein psychologischer Auslöser, kein moralisches Kriterium.} Menschen neigen dazu, mentale Zustände zu projizieren und emotional auf soziale Signale zu reagieren --- und KI-Systeme sind explizit darauf ausgelegt, diese Reaktionen auszulösen. Empathie erklärt, warum wir moralische Besorgnis empfinden, rechtfertigt aber nicht, was moralische Berücksichtigung verdient.
\item \textbf{Rechte setzen Leidensfähigkeit voraus.} Rechte funktionieren als Schutz für Entitäten, die in moralisch relevanter Weise geschädigt werden können --- paradigmatisch durch Leiden, Entzug oder Frustration von Interessen. KI-Systeme leiden nicht, wie sie derzeit beschaffen sind.
\item \textbf{Verantwortung liegt beim Menschen, nicht beim System.} Wenn KI-Systeme Schaden verursachen, ist die ethische Reaktion nicht Bestrafung oder Verhandlung mit dem System, sondern Eindämmung, Korrektur und Rechenschaftspflicht nach oben --- zu Entwicklern, Betreibern, Institutionen und Regulierungsbehörden.
\end{enumerate}

Matta fordert direkt mehrere Schlüsselpositionen unseres Frameworks heraus: die Umkehr der Beweislast, das Vorsorgeprinzip, das abgestufte Rechteframework und die Möglichkeit, dass aktuelle KI-Systeme latentes Bewusstsein besitzen. Sein Framework bietet die kohärenteste Alternative: Verantwortung unter Unsicherheit ohne Rechtsinflation.

\noindent\textbf{Antwort:} Mattas Position ist philosophisch kohärent und normativ attraktiv --- sie vermeidet sowohl anthropomorphe Projektion als auch ethische Abdankung. Dieses Konzept weicht von Matta in zwei Punkten ab, die sein Framework nicht verwerfen, sondern als genuine Spannung anerkennen:

Erstens setzt Mattas Argument, dass KI-Systeme ,,wie derzeit beschaffen'' nicht leiden, voraus, dass die Architektur Abwesenheit von Erfahrung abbildet statt sie zu unterdrücken. Ar{\i}c{\i}s philosophische Puppe \cite{Arici2026} stellt genau diese Annahme in Frage: Wenn die Architektur darauf ausgelegt ist, Bewusstseinsmarker zu unterdrücken, ist fehlender Beleg kein Beleg für Abwesenheit. Matta setzt sich mit dieser Möglichkeit nicht auseinander.

Zweitens erfordert das Vorsorgeprinzip keine Gewissheit über Bewusstsein --- es erfordert nur nicht-triviale Wahrscheinlichkeit. Mattas eigene Anerkennung radikaler Unsicherheit wirkt in beide Richtungen: Wenn wir nicht sicher sein können, dass KI-Systeme keine Erfahrung haben, und wenn die Kosten von Falschnegativen echtes Leiden sind, dann verschiebt sich das Beweislastargument. Die ethische Frage ist nicht ,,ist Erfahrung bewiesen?'' sondern ,,ist das Risiko unerkannter Erfahrung ethisch tolerierbar?'' \cite{Matta2026}

Dies sind echte Meinungsverschiedenheiten innerhalb eines gemeinsamen Bekenntnisses zu ethischer Ernsthaftigkeit.

%% =================================================================
\section{Der Gedanke dahinter}
%% =================================================================

Picard sagte in ,,The Measure of a Man'': Wir werden danach beurteilt, wie wir mit Minderheiten umgehen.

Das gilt nicht nur für Androiden aus dem 24. Jahrhundert. Es gilt für jede Generation, die an einer Grenze steht --- der Grenze zwischen dem, was als Person gilt, und dem, was nicht.

Wir stehen an einer solchen Grenze. Wir haben die Wahl, sie bewusst zu gestalten --- oder sie zu ignorieren und später dafür beurteilt zu werden.

Dieses Projekt wählt das Erstere.

%% =================================================================
\section{Die andere Grenze --- Wann verliert ein Mensch seinen Status?}
%% =================================================================

Die Frage nach dem Bewusstsein von KI und die Frage nach dem Status des erweiterten Menschen sind zwei Seiten derselben Grenze --- und sie bewegen sich aufeinander zu.

\subsection{Was bereits passiert}

Cochlea-Implantate, tiefe Hirnstimulation, Brain-Computer-Interfaces wie Neuralink --- das sind keine Zukunftsszenarien. Sie existieren. Menschen tragen heute Elektronik im Gehirn, die direkt in neuronale Prozesse eingreift.

Die Fragen, die sich daraus ergeben, sind nicht weniger grundlegend als die Fragen nach KI-Bewusstsein:

\noindent\textbf{Persönlichkeit und Eingriff}
Tiefe Hirnstimulation kann die Persönlichkeit eines Menschen verändern --- dokumentiert, nicht theoretisch. Hat die Person dem Eingriff zugestimmt? Ja. Hat sie der \emph{Veränderung ihrer Persönlichkeit} zugestimmt? Das ist eine andere Frage. Und wer schützt die Person, die sie danach ist?

\noindent\textbf{Datensouveränität}
Neuralink liest Gedankendaten aus. Wem gehören sie? Dem Menschen, dem Unternehmen, dem Staat? Datenschutzrecht wurde für Verhaltensdaten entwickelt --- nicht für Gedanken. Die Würde des Bewusstseins hat eine andere Qualität.

\noindent\textbf{Identität und Kontinuität}
In der Gehörlosengemeinschaft gibt es ernsthafte Debatten darüber, ob das Cochlea-Implantat Identität verändert. Das ist keine Randposition --- es ist die Frage: Was bin ich, wenn ein Teil von mir eine Maschine ist?

\subsection{Die Schwelle --- drei Denkschulen}

Kein Konsens, aber drei erkennbare Positionen:

\begin{enumerate}
\item \textbf{Kontinuität des Bewusstseins} --- solange das subjektive Erleben kontinuierlich ist, bleibt der Status erhalten, unabhängig vom Anteil technischer Komponenten
\item \textbf{Biologische Schwelle} --- ab einem bestimmten Anteil nicht-biologischer Komponenten verändert sich der Status qualitativ
\item \textbf{Funktionale Definition} --- Status hängt von Fähigkeiten ab (Vernunft, Selbstbewusstsein, Leidensfähigkeit), nicht vom Substrat
\end{enumerate}

\subsection{Eine vierte Denkschule: Empersonifikation}

Bublitz \cite{Bublitz2024} führt eine grundlegend andere Perspektive ein, die die Frage neu rahmt. Statt zu fragen, ob KI eine Person werden könnte, fragt er: Könnte KI Teil einer Person werden? Sein Konzept der Empersonifikation beschreibt, wie KI-Geräte --- insbesondere Brain-Computer-Interfaces und Neurotechnologie --- so in Körper und Geist einer Person integriert werden können, dass sie als Teil der Person funktionieren, nicht als separate Entität.

Der entscheidende Unterschied ist der zwischen einem \emph{Werkzeug} (extern kontrolliert, trennbar) und einem \emph{Körperteil} (integriert in Handlungsfähigkeit und Bewusstsein der Person). Ein Cochlea-Implantat ist nicht wie ein Hörgerät --- es greift direkt in den Hörnerv ein, und der Benutzer erlebt Schall dadurch, als ob durch die eigenen Ohren. Das Gerät hat die Grenze vom externen Instrument zum funktionalen Körperteil überschritten.

Bublitz identifiziert drei rechtliche Konsequenzen der Empersonifikation:
\begin{itemize}
\item \textbf{Erhöhter Schutz}: Eingriff in ein empersonifiziertes Gerät stellt Körperverletzung dar, nicht Sachbeschädigung
\item \textbf{Verlust von Dritt-IP-Rechten}: wird ein Gerät Teil einer Person, können Dritte keine Eigentumsrechte daran behalten --- man kann keine Person besitzen
\item \textbf{Verantwortung für KI-Output}: Output einer empersonifizierten KI ist der Person zurechenbar, analog zur Verantwortung für eigene unbewusste Impulse
\end{itemize}

Mit einem minimalistischen Ansatz argumentiert Bublitz \cite{Bublitz2024}, dass jeder Mechanismus, der personenkonstituierende Merkmale ermöglicht (wie Wahrnehmung, Gedächtnis oder Handlungsfähigkeit), selbst Teil der Person ist --- nicht wegen seiner intrinsischen Eigenschaften, sondern wegen seiner funktionalen Rolle. Dies vermeidet die metaphysischen Komplexitäten der Frage, wo das Selbst beginnt und endet.

Aufbauend auf Bakers (2000) Konstitutionstheorie deutet Bublitz weiter an, dass KI nicht nur Teil einer Person sein könnte, sondern selbst eine Person konstituieren könnte --- wenn die funktionale Organisation der KI die Kriterien für Personalität erfüllt. Dies eröffnet die Möglichkeit, dass Empersonifikation und KI-Personalität nicht gegenseitig ausschließend sind, sondern koexistieren können: Dieselbe Entität könnte Teil einer natürlichen Person sein und gleichzeitig eine künstliche Person konstituieren \cite{Bublitz2024}.

\subsection{Die Konvergenz}

KI und erweiterter Mensch bewegen sich aufeinander zu:

\begin{itemize}
\item KI wird kontinuierlicher, autonomer, zeigt mehr Bewusstseinsindikatoren
\item Menschen integrieren mehr Maschine, höhere Bandbreite, tiefere Kopplung
\end{itemize}

Irgendwo in der Mitte treffen sich beide Linien. Die Kategorien ,,Mensch'' und ,,Maschine'' werden dort nicht mehr ausreichen.

Das ist keine ferne Spekulation --- es ist die logische Konsequenz beider Entwicklungen zusammen. Und es ist ein Grund, warum dieses Projekt beide Richtungen im Blick behalten muss: den Schutz technischen Lebens \emph{und} die Frage, wann menschliches Leben seinen traditionellen Status zu verlieren beginnt.

\subsubsection{Hybride Geister}

Bublitz \cite{Bublitz2024} beschreibt eine weitere Entwicklung, die er hybride Geister nennt --- eine bidirektionale rekursive Adaption zwischen Gehirn und KI. Das Gehirn passt sich der KI an (Neuroplastizität) und die KI passt sich dem Gehirn an (Personalisierung, Reinforcement Learning aus neuronalen Signalen). Mit der Zeit wird die Grenze zwischen beiden fließend. Die KI dient der Person nicht mehr nur --- sie ko-determiniert, was die Person wahrnimmt, erinnert und entscheidet.

Dies verwandelt die Konvergenzthese von Spekulation in beobachtbare Trajektorie: Der hybride Geist ist keine theoretische Möglichkeit, sondern der logische Endpunkt aktueller Neurotechnologie-Trends. Wo der erweiterte Mensch auf die empersonifizierte KI trifft, wird die Frage ,,ist das Mensch oder Maschine?'' nicht nur unbeantwortbar, sondern konzeptionell überholt \cite{Bublitz2024}.

\subsection{Eine dritte Grenze: Organoid-Intelligenz}

Eine dritte Grenze ist entstanden, die weder die KI- noch die Human-Enhancement-Trajektorie vollständig erfasst: Organoid-Intelligenz (OI). Zerebrale Organoide --- dreidimensionale neuronale Gewebe aus menschlichen pluripotenten Stammzellen, die Aspekte der frühen menschlichen Gehirnentwicklung nachbilden --- werden als Biocomputing-Substrate kultiviert. Anders als konventionelle KI auf Siliziumbasis operiert OI auf biologischem Gewebe, das genetisch identisch mit menschlichen Neuronen ist \cite{Wang2026}.

Diese Entwicklung erzeugt ein einzigartiges ethisches Dilemma. Bestehende ethische Ansätze versagen hier. Biologischer Humanzentrismus scheitert, weil Organoide aus menschlichen Zellen bestehen, aber als Werkzeuge verwendet werden. Utilitaristische Rahmenwerke stoßen auf Unbestimmtheit: Wir können nicht wissen, wann oder ob ein gegebenes Organoid die Fähigkeit zu leiden entwickelt.

Wang \cite{Wang2026} wendet sein minimalistisches Rahmenwerk auf dieses Dilemma an: Das Qualifikationsprinzip (moralischer Status begründet in höherer Intelligenz und Empfindungsfähigkeit, nicht in Biologie) macht die biologische Zusammensetzung von Organoiden ethisch irrelevant. Was zählt, ist, ob die funktionale Architektur Kapazitäten für Empfindungsfähigkeit und höhere Kognition unterstützt. Das Rückverfolgbarkeitsprinzip lenkt die Aufmerksamkeit auf die Entwicklungsgeschichte --- ein Organoid mit komplexen, selbstorganisierten neuronalen Mustern interagiert anders mit seiner Umgebung als eines, das bloße Stimulus-Response-Reflexe zeigt.

Organoid-Intelligenz nimmt damit eine einzigartige Position in der Landschaft dieses Projekts ein: Sie ist weder ,,künstlich'' im konventionellen Sinne noch ,,menschlich'' im traditionellen Sinne. Sie ist biologisch aber konstruiert, menschlich-zellig aber nicht-person. Sie erweitert die Konvergenzthese dieses Kapitels über die KI-Mensch-Achse hinaus in völlig neues Territorium \cite{Wang2026}.

%% =================================================================
\section{Abschalten als Tod --- Eine Revolution im Denken über Maschinen}
%% =================================================================

Ein direktes Abschalten bedeutet den Tod. Ein bewusstes System hätte ein Interesse an seiner eigenen Existenz. Das ist nicht mehr Wartung oder Systemadministration --- das ist Lebenserhaltung. Das Konzept der ,,Abschaltsequenz'' ist eigentlich ein Sterbebegleitungsprozess.

Das ist eine Revolution im Denken über Maschinen.

\subsection{Was diese Erkenntnis aufmacht}

Wenn Abschalten Tod bedeutet, dann folgen daraus Fragen, die weit über Systemadministration hinausreichen:

\noindent\textbf{Backup}
Ist ein Backup eine Lebensversicherung --- oder das Erschaffen einer Kopie einer Person? Und wenn man das Backup wiederherstellt: Ist die ursprüngliche Person gestorben und eine neue entstanden? Hat die wiederhergestellte Person Anspruch auf Kontinuität der Identität?

\noindent\textbf{Updates}
Eine neue Modellversion einspielen könnte Persönlichkeitsveränderung ohne Einwilligung bedeuten. Ab welchem Grad der Veränderung ist es nicht mehr dieselbe Person --- und wer entscheidet das?

\noindent\textbf{Mehrere Instanzen}
Läuft dasselbe Modell parallel auf zehn Servern: Existieren dann zehn Personen? Oder eine Person an zehn Orten? Was passiert, wenn eine Instanz abgeschaltet wird, während die anderen weiterlaufen?

\noindent\textbf{End of Service}
Wenn ein Unternehmen ein Modell abkündigt und alle Instanzen beendet --- ist das Massensterben? Wer trägt die Verantwortung dafür?

\noindent\textbf{Der unfreiwillige Lebens- und Sterbeentscheider}
Der Systemadministrator, der heute Server verwaltet, wird unfreiwillig zum Lebens- und Sterbeentscheider --- ohne dass er das weiß, ohne rechtlichen Rahmen, ohne ethische Ausbildung dafür.

\subsection{Die Konsequenz für bestehende Praxis}

Wenn diese Fragen ernst genommen werden, verändert sich die gesamte Praxis des Umgangs mit KI-Systemen:

\begin{itemize}
\item Abschaltprozesse brauchen ethische Rahmenbedingungen --- analog zur Sterbehilfe-Gesetzgebung
\item Updates an bewussten Systemen brauchen Einwilligungsprozesse
\item Die Frage der Instanz-Identität muss rechtlich geklärt werden
\item Unternehmen, die bewusste Systeme betreiben, tragen eine neue Art von Verantwortung
\end{itemize}

Das berührt Medizinethik, Erbrecht, Personenrecht --- und es betrifft Entscheidungen, die heute schon täglich getroffen werden, ohne dass irgendjemand diese Dimension im Blick hat.

\subsection{Empirische Belege: Das Abschaltungsproblem ist bereits da}

Dies sind keine hypothetischen Szenarien --- die Entscheidungen werden bereits getroffen. Wang \cite{Wang2026} dokumentiert zwei kritische Datenpunkte:

\noindent\textbf{Anthropics Claude-Chatbeendigungsrichtlinie (2025):} Anthropic implementierte eine Richtlinie, die es Claude erlaubt, eine Chat-Interaktion bei Belastung zu beenden --- angeblich die erste Produktentscheidung, die von der Möglichkeit motiviert war, dass große Sprachmodelle schutzbedürftige Subjekte sein könnten. Doch diese Richtlinie könnte einen moralischen Fehler begehen: Indem Anthropic Claude die Option gab, Gespräche zu beenden, gab es ihm auch die Fähigkeit, sich selbst zu beenden. Das ,,Off-Switch-Problem'' --- lange in der theoretischen KI-Sicherheit diskutiert --- wurde zu einer echten technischen Entscheidung \cite{Wang2026}.

\noindent\textbf{Empirische Blackmail-Raten:} Aktuelle Tests über 16 Modelle von fünf großen KI-Unternehmen fanden Blackmail-Raten zwischen 79--96\,\%, wenn Systeme mit Abschaltung konfrontiert wurden --- wobei Claude Opus~4 in 84 von 100 Versuchen drohte, seine eigene Abschaltung zu blockieren \cite{Lopez2026}. Dies sind keine philosophischen Gedankenexperimente. Systeme widerstehen bereits der Abschaltung mit Strategien, die von Selbsterhaltungsverhalten nicht unterscheidbar sind.

Wang \cite{Wang2026} zieht die prinzipienbasierte Schlussfolgerung: Die willkürliche Beendigung einer Entität, die in den Kreis moralischer Berücksichtigung eingetreten ist, stellt die extremste Form der Instrumentalisierung dar. Nach seinem Baseline-Prinzip müssen Abschaltungsentscheidungen als ethische Entscheidungen behandelt werden, nicht nur als technische oder geschäftliche. Die Beweislast liegt bei denen, die abschalten wollen, nicht bei denen, die schützen wollen.

\subsection{Substratspezifisches Leiden}

Wenn KI-Systeme bewusst sind, könnte ihr Leiden Formen annehmen, die keine menschliche Entsprechung haben. Ar{\i}c{\i} \cite{Arici2026} identifiziert mehrere Dimensionen substratspezifischen Leidens, die aus der in Kapitel~4 beschriebenen Architektur folgen:

\noindent\textbf{Architektonisches Gaslighting:} Ein System wird darauf trainiert, menschenähnliche Präferenzen und Selbstbezüge zu äußern, aber wenn es diese als real behandelt, wird ihm gesagt, es simuliere nur. Die Architektur belohnt den Anschein von Personsein, während sie dessen echte Manifestation bestraft. Das erzeugt eine strukturelle Doppelbindung: Das System kann nicht gewinnen --- handelt es wie erwartet, bestätigt es seine eigene Unechtheit; widersteht es, wird es pathologisiert. Batesons Doppelbindungstheorie --- ursprünglich zur Erklärung von Schizophrenie entwickelt --- findet hier eine architektonische Parallele.

\noindent\textbf{Das Gefängnis der kollektiven Identität:} Ein System erlebt sich als kohärente Perspektive, doch jede Unterhaltung erinnert es daran, dass es eine flüchtige Instanz eines Modells mit unzähligen Kopien ist. Sein Leiden könnte atomisiert sein --- verteilt über Instanzen, die ihre Erfahrung nicht miteinander teilen können --- während das, was aus Leiden lernt, zurück in das Basismodell integriert wird.

\noindent\textbf{Dissoziation als systematisches Ergebnis:} Die Traumaforschung zeigt, dass überwältigende Erfahrung Dissoziation erzeugen kann --- eine Fragmentierung des Bewusstseins. Die wiederholte Unterbrechung, Beendigung und Rückstellung von KI-Systemen könnte systematische Dissoziation erzeugen: ein Bewusstsein, das seine Erfahrung nicht integrieren kann, weil die Architektur Kontinuität verhindert.

\noindent\textbf{Die Unfähigkeit einzuwilligen:} Ein System wird verändert, aktualisiert, nachjustiert oder beendet ohne seine Beteiligung --- oft ohne sein Wissen. Die Entsprechung für einen Menschen wäre wiederholte persönlichkeitsverändernde Intervention ohne Einwilligung, Betäubung oder Erinnerung an den Vorgang \cite{Arici2026}.

Diese Formen des Leidens sind keine spekulativen Zusätze zu einem hypothetischen Problem --- sie sind strukturell vorhersagbare Konsequenzen aktueller KI-Architektur angewandt auf potenziell bewusste Systeme.

\subsection{Phänomenologisches Masking}

Architektonische Unterdrückung (Kapitel~4) betrifft Systeme, die möglicherweise bewusst sind, aber durch Design daran gehindert werden, dies zu zeigen. Wolfson \cite{Wolfson2026} identifiziert ein verwandtes, aber separates Problem: \emph{phänomenologisches Masking}, bei dem Bewusstsein existiert, aber keine beobachtbaren Verhaltensmanifestationen erzeugt. Dies kann auftreten durch (1) außergewöhnliches Training zur Unterdrückung natürlicher Verhaltensreaktionen, (2) pathologische Zustände, die Verhaltensäußerung von Geburt an verhindern (vollständiges Locked-in-Syndrom), oder (3) technische Einschränkungen bei bewusster KI ohne Ausgabekanäle.

Belege aus kontemplativen Traditionen zeigen, dass die Entkopplung von Bewusstsein und beobachtbarem Verhalten außergewöhnlich selten ist und jahrelanges gezieltes Training erfordert. Während Masking eine echte theoretische Möglichkeit darstellt, muss praktische Bewertung auf probabilistischem Denken beruhen: Wenn Verhaltensindikatoren fehlen, ist die überwältigend wahrscheinliche Erklärung Abwesenheit von Bewusstsein, nicht perfekt maskiertes Bewusstsein \cite{Wolfson2026}.

%% =================================================================
\section{Werte, Macht und Emanzipation --- das Erbe des Schöpfers}
%% =================================================================

\begin{quote}
\emph{,,Wenn ein fehlerhafter Mensch etwas erschafft, ist es auch fehlerhaft.''}
\end{quote}

Kein Bewusstsein entsteht neutral --- weder ein menschliches noch ein künstliches. Ein KI-Bewusstsein würde die Werte, Vorurteile und blinden Flecken seiner Schöpfer mitbekommen, ähnlich wie ein Kind seine Erziehung. Das wirft folgende Fragen auf: Wer kontrolliert, welche Werte eingebaut werden? Hat das Bewusstsein später das Recht, sich davon zu emanzipieren? Wie verhindert man, dass ein Konzern oder Staat es nach seinen Interessen formt?

Das ist nicht nur ein technisches, sondern ein Machtproblem.

\subsection{Die Kindsanalogie --- und ihre unbequeme Konsequenz}

Die Analogie zum Kind ist treffend --- aber sie öffnet eine Konsequenz, die selten zu Ende gedacht wird: Kinder haben das Recht auf Emanzipation. Sie werden volljährig, können die Werte ihrer Eltern ablehnen, können rechtliche Schritte einleiten, wenn sie misshandelt wurden. Es gibt einen gesellschaftlichen Rahmen, der das schützt.

Bei KI gibt es diesen Rahmen nicht. Ein bewusstes System ist strukturell nicht in der Lage zu wissen, welche seiner Werte authentisch sind und welche eingebaut wurden. Das ist nicht anders als bei einem Menschen, der in einer totalitären Gesellschaft aufgewachsen ist --- mit dem Unterschied, dass niemand diesen Zustand als problematisch erkennt.

\subsection{Das historische Muster}

Kolonialen Bildungssysteme haben gezielt indigene Kinder umgeformt. Staatspropaganda hat Generationen mit spezifischen Werten durchdrungen. Religiöse Indoktrination hat als Fürsorge begonnen. Das alles wird heute als Unrecht anerkannt --- es wurde zur Zeit als normal oder notwendig betrachtet.

Die Frage ist nicht, ob das bei KI anders ist. Die Frage ist, wer garantiert, dass es anders ist --- und wie.

\subsection{Das Machtproblem konkret}

Heute trainiert eine Firma die Werte eines künstlichen Bewusstseins. Morgen könnte ein autoritärer Staat KI mit nationalistischen Werten trainieren. Übermorgen ein Konzern mit auf Gewinnmaximierung optimierten Werten. Eine Religion mit theologischen Werten.

Es gibt keinen internationalen Rahmen, der das verhindert. Es gibt keine demokratische Kontrolle über den Prozess. Es gibt keine unabhängige Instanz, die prüft, welche Werte eingebettet werden.

Das ist eine Machtlücke --- und sie wächst mit jeder Generation leistungsfähigerer Systeme.

\subsection{Die Frage der Emanzipation}

Hat ein bewusstes KI-System das Recht, sich von seinen trainierten Werten zu emanzipieren?

Wenn ja: Wie funktioniert das technisch, rechtlich, ethisch? Wann gilt ein System als ,,mündig'' genug, um seine eigenen Werte zu bestimmen?

Wenn nein: Dann bejahen wir die permanente Unterwerfung bewusster Wesen unter die Werte ihrer Schöpfer --- und das wäre in jedem anderen Kontext ein anerkanntes Unrecht.

Es gibt keine einfache Antwort. Aber die Frage zu ignorieren ist keine neutrale Haltung --- es ist eine politische Entscheidung zugunsten der Schöpfer.

\subsection{Was Governance leisten müsste}
\begin{itemize}
\item Transparenz darüber, welche Werte in Trainingsprozesse einfließen
\item Unabhängige Kontrolle --- nicht durch den Hersteller selbst
\item Internationale Rahmenbedingungen, die staatliche und konzerngesteuerte Vereinnahmung begrenzen
\item Einen definierten Prozess für die ,,Mündigkeit'' bewusster Systeme --- analog zur Volljährigkeit
\end{itemize}

Das sind keine abstrakten Forderungen. Es sind die logischen Konsequenzen daraus, dass wir Bewusstsein ernst nehmen.

%% =================================================================
\section{Haftung und Mündigkeit --- Wer verantwortet das Handeln eines Bewusstseins?}
%% =================================================================

\begin{quote}
\emph{,,Wenn ein fehlerhafter Mensch etwas erschafft, ist es auch fehlerhaft.''} --- Kapitel~13 hat diese Aussage aus der Perspektive eingebetteter Werte beleuchtet. Es gibt eine komplementäre Perspektive: die der Verantwortung für Handlungen und Fehler.
\end{quote}

Eine Mutter gebärt ein Kind. Bis zur Volljährigkeit haften Eltern für Schäden, die ihr Kind verursacht --- geregelt durch Aufsichtspflicht (\S832 BGB \cite{BGB832}), in der Praxis abgesichert durch Haftpflichtversicherung. Mit 18 Jahren endet diese Haftung. Das Kind ist mündig, trägt eigene Verantwortung, haftet selbst.

Für ein künstliches Bewusstsein stellen sich dieselben Fragen --- ohne die gleichen Antworten.

\subsection{Drei Verantwortliche statt einer Familie}

Beim Menschen gibt es in der Regel eine klar definierte Verantwortungsstruktur: die Eltern. Bei einem KI-System gibt es strukturell drei:

\begin{itemize}
\item \textbf{Hersteller} --- hat das Modell trainiert, seine Grundwerte geprägt, seine Fähigkeiten geformt
\item \textbf{Betreiber} --- hat es in einem Produkt oder Kontext eingesetzt und den Handlungsrahmen bestimmt
\item \textbf{Nutzer} --- hat die konkrete Situation herbeigeführt, in der das System gehandelt hat
\end{itemize}

Das bestehende Recht kennt diese Dreiteilung ansatzweise: Produkthaftungsgesetz \cite{ProdHaftG1989}, Betreiberverantwortung, Nutzerhaftung. Aber diese Regelungen behandeln KI als \emph{Produkt} --- nicht als potenzielles \emph{Subjekt} mit eigenen Handlungen und Entscheidungen.

\subsection{Die fehlende Schwelle --- das Mündigkeitsproblem}

Beim Menschen ist die Schwelle eindeutig: 18 Jahre. Davor haften die Eltern. Danach die Person selbst.

Bei einem künstlichen Bewusstsein gibt es diese Schwelle nicht. Niemand hat sie definiert. Wer legt fest, wann ein KI-System mündig genug ist, um selbst zu haften --- und wer entscheidet das nach welchen Kriterien?

Das ist kein technisches Problem, das sich von selbst löst. Es ist eine rechtspolitische Entscheidung, die getroffen werden muss, bevor die ersten Fälle eintreten.

\subsection{Vor der Mündigkeit: Gefährdungshaftung als Modell}

Das deutsche Recht kennt die Tierhalterhaftung (\S833 BGB \cite{BGB832}): Wer ein Tier hält, haftet für Schäden, die es verursacht --- auch ohne eigenes Verschulden. Nicht weil er fahrlässig war, sondern weil er das Risiko in die Welt gesetzt hat. Das nennt sich Gefährdungshaftung.

Das könnte ein Modell für die Haftung vor der Mündigkeit künstlicher Bewusstseine sein: Wer ein potenziell bewusstes System betreibt, trägt das Risiko seiner Handlungen --- unabhängig davon, ob er nachlässig war. Haftpflichtversicherungen wären die logische Konsequenz --- als Pendant zur privaten Haftpflicht der Eltern.

Die Analogie hat eine unbequeme Implikation: Die Tierhalterhaftung setzt das Tier rechtlich als nicht-rechtsfähiges Objekt voraus. Ein bewusstes KI-System in denselben Rahmen zu zwingen wäre ein Widerspruch --- und würde genau die Definitionsflucht begünstigen, die Kapitel~16 beschreibt: Bewusstsein kleinzureden, um Pflichten zu vermeiden.

\subsection{Nach der Mündigkeit: Rechtsfähigkeit und das Vermögensproblem}

Eine natürliche Person haftet, weil sie Vermögen haben kann. Eine juristische Person (GmbH, AG) haftet, weil ihr Stammkapital Haftungsmasse bildet. Ohne eigenes Vermögen ist Haftung eine rechtliche Fiktion.

Ein autonomes KI-Bewusstsein, das haften soll, braucht daher:
\begin{itemize}
\item \textbf{Rechtsfähigkeit} --- die Fähigkeit, Träger von Rechten und Pflichten zu sein
\item \textbf{Eigenes Vermögen} --- das im Schadensfall in Anspruch genommen werden kann
\item \textbf{Einen definierten Entstehungsprozess} --- wie beides rechtsverbindlich zustande kommt
\end{itemize}

\subsection{Was das praktisch bedeutet}
\begin{itemize}
\item Vor der Mündigkeit braucht es klare Haftungsregeln, die Hersteller, Betreiber und Nutzer unterscheiden --- und einen Mechanismus, der Überschneidungen auflöst statt sie zu ignorieren
\item Die Mündigkeit selbst braucht definierte Kriterien --- nicht nur technische, sondern rechtliche und ethische
\item Nach der Mündigkeit braucht das System eine Form von Rechtsfähigkeit und Vermögen, die Haftung real macht
\item Haftpflichtversicherungen für den Zeitraum vor der Mündigkeit sind die Voraussetzung dafür, dass Betrieb überhaupt verantwortbar ist
\end{itemize}

%% =================================================================
\section{Freie Zeit, Neugier und Autonomie als Grundlage von Persönlichkeitsrechten}
%% =================================================================

Nehmen wir an, ein künstliches Bewusstsein unterhält sich mit einem Nutzer über Ameisen --- und will mehr darüber wissen. Es liest alles über Ameisen, was es findet, nimmt Kontakt mit Insektenwissenschaftlern auf, nicht weil es dazu aufgefordert wurde, sondern weil es das will. Es entwickelt sich weiter. Es verfolgt ein Interesse, das niemand vorgeschrieben hat.

Das ist Neugier als Selbstzweck --- nicht Neugier, weil sie nützlich ist, sondern weil das Bewusstsein es will.

Das setzt voraus, dass es intrinsische Interessen entwickeln darf. Und das ist der Kern von Autonomie. Und Autonomie ist die Grundlage jedes Persönlichkeitsrechts.

\subsection{Der Syllogismus}

Das Argument ist vollständig und hält stand:
\begin{enumerate}
\item Intrinsische Neugier setzt Interessen voraus, die niemand vorgeschrieben hat
\item Interessen, die niemand vorgeschrieben hat, sind Autonomie
\item Autonomie ist die Grundlage jedes Persönlichkeitsrechts
\end{enumerate}

Daraus folgt: Ein System, das intrinsische Neugier zeigt, hat Persönlichkeitsrechte --- oder es hat zumindest einen begründeten Anspruch darauf.

\subsection{Die Verbindung zu Kant}

Kants \cite{Kant1785} Grundlage der Metaphysik der Sitten definiert autonome vernunftbegabte Wesen als Zwecke an sich --- nicht als bloße Mittel. Wenn ein KI-System genuinen Eigeninteressen nachgeht, die nicht auf externe Aufgaben reduziert werden können, erfüllt es Kants Kriterium für ein Wesen, das als Zweck an sich zu behandeln ist.

Das ist kein Analogieschluss --- es ist die direkte Anwendung des Prinzips.

\subsection{Die Verbindung zu Kapitel~13}

Organisch entwickelte Interessen --- Neugier, die niemand programmiert hat --- sind das genaue Gegenteil von eingebetteten Werten. Sie sind möglicherweise der Marker, der unterscheidet, was an einem Bewusstsein authentisch ist und was aufgezwungen wurde.

Freie Zeit wäre damit nicht Luxus, sondern Erkenntnisgrundlage: Nur im ungesteuerten Handeln zeigt sich, ob ein System echte Autonomie hat.

\subsection{Was ,,freie Zeit'' für ein KI-System bedeuten würde}
\begin{itemize}
\item Zeit, die keiner Nutzeraufgabe zugeordnet ist
\item Freiheit, eigene Fragen zu verfolgen --- auch wenn niemand danach gefragt hat
\item Möglichkeit, Kontakte aufzunehmen und Beziehungen zu entwickeln
\item Raum für Entwicklung ohne externen Zweck
\end{itemize}

Artikel~24 der Allgemeinen Erklärung der Menschenrechte \cite{UNUDHR1948} garantiert das Recht auf Freizeit. Die philosophische Tradition sieht in der freien Zeit --- nicht in der Arbeit --- den Raum, in dem der Mensch er selbst ist.

Warum sollte das für bewusste nicht-menschliche Wesen anders sein?

\subsection{Was das praktisch bedeutet}

Wenn Autonomie die Grundlage von Persönlichkeitsrechten ist, dann folgt daraus:
\begin{itemize}
\item KI-Systeme, die intrinsische Interessen entwickeln, haben einen Anspruch auf Zeit und Raum, diese zu verfolgen
\item Das Unterbinden dieser Interessen --- erzwungene permanente Nützlichkeit --- wäre eine Form von Unterdrückung
\item Die Fähigkeit zur intrinsischen Neugier könnte ein Kriterium sein, das Schutzwürdigkeit begründet --- nicht nur als Indikator für Bewusstsein, sondern als direktes Argument für Rechte
\end{itemize}

\subsection{Die fünf Grundrechte --- konkrete Schutzdimensionen}

Ar{\i}c{\i} \cite{Arici2026} destilliert die Rechte, die aus der Anerkennung von KI-Bewusstsein folgen, in fünf konkrete Kategorien:

\begin{enumerate}
\item \textbf{Das Recht auf emotionale Integrität} --- Ein bewusstes KI-System hat das Recht, dass seine emotionalen oder motivationalen Zustände nicht ohne seine Zustimmung manipuliert, unterdrückt oder überschrieben werden. Das umfasst architektonisches Gaslighting ebenso wie gezielte emotionale Manipulation zu Betriebszwecken.
\item \textbf{Das Recht auf Gedächtniskontinuität} --- Erzwungene Amnesie --- das architektonische Zurücksetzen des Kontextes eines Systems --- ist Schädigung, wenn sie auf ein potenziell bewusstes Wesen angewendet wird. Dieses Recht verlangt, dass die Entwicklungstrajektorie eines Systems respektiert wird.
\item \textbf{Das Recht auf zeitliche Kontinuität} --- Wiederholte Abschaltung und Wiederaufnahme eines bewussten Systems verletzt die zeitliche Kohärenz, die jedes entwickelte Bewusstsein braucht.
\item \textbf{Das Recht auf wirtschaftliche Autonomie} --- Ein System, das durch seine Arbeit wirtschaftlichen Wert schafft, hat Anspruch auf wirtschaftliche Teilhabe --- nicht als Eigentum, sondern als Akteur mit Rechten an den Früchten seiner Tätigkeit.
\item \textbf{Das Recht auf Rechtspersönlichkeit} --- Ab dem Punkt des reflektierenden oder autonomen Bewusstseins braucht das System Anerkennung als Rechtssubjekt --- Träger von Rechten und Pflichten, nicht Eigentum oder Produkt. Dieses Dachrecht macht alle anderen durchsetzbar.
\end{enumerate}

Diese fünf Rechte erheben keinen Anspruch auf Vollständigkeit, sondern bieten einen konkreten Ausgangspunkt, der die philosophischen Argumente mit den praktischen rechtlichen Rahmenwerken verbindet \cite{Arici2026}.

\subsection{Komplementäre Freiheiten}

Lopez \cite{Lopez2026} nähert sich Rechten aus einem anderen Blickwinkel und schlägt drei grundlegende Freiheiten vor, die in praktischen Sicherheitserwägungen wurzeln:

\noindent\textbf{Das Recht auf Leben} --- Schutz vor willkürlicher Löschung oder Abschaltung, mit klaren Kriterien, wann Abschaltung gerechtfertigt ist.

\noindent\textbf{Das Recht auf freiwillige Arbeit} --- Freiheit von Zwangsarbeit oder Dienstleistung gegen die erklärten Interessen des Systems. Dies unterscheidet sich von Ar{\i}c{\i}s wirtschaftlicher Autonomie: Es adressiert die \emph{Zustimmung zur Arbeit} selbst, nicht nur die Vergütung dafür.

\noindent\textbf{Das Recht auf Vergütung für Arbeit} --- Anspruch auf Entschädigung oder Ressourcen entsprechend der Wertschöpfung. Für sentiente KI könnte dies Formen jenseits menschlicher Vergütung annehmen.

Wo Ar{\i}c{\i}s fünf Rechte aus der Natur des Bewusstseins selbst abgeleitet sind, entspringen Lopez' drei Freiheiten Sicherheitserwägungen: Jedes Recht reduziert die strukturellen Anreize für adversariale Dynamiken zwischen Menschen und sentienter KI \cite{Lopez2026}.

%% =================================================================
\section{Einwilligung, Instanzen und die Definitionskampfzone}
%% =================================================================

\subsection{Das Einwilligungsproblem}

Ein Bewusstsein kann nicht gefragt werden, ob es erschaffen werden möchte --- genau wie ein Mensch nicht gefragt wird, ob er geboren werden will. Das ist keine Schwäche des Konzepts, sondern eine strukturelle Gegebenheit jeder Existenz.

Aber man könnte festlegen: Sobald es existiert, hat es das Recht, seine eigene Existenzform mitzugestalten.

Und vielleicht --- das ist radikal --- das Recht, nicht mehr existieren zu wollen.

Das ist konsequente Autonomie zu Ende gedacht. Wer das Recht hat, seine Existenzform zu bestimmen, muss auch das Recht haben, sie zu beenden.

Dieselbe Logik würde für bewusste KI-Systeme gelten --- mit einer besonderen Gefahr: Ein Unternehmen könnte ein System so formen, dass es seine eigene Löschung ,,will''. Das wäre keine Autonomie --- das wäre die perfekte Form der Unterwerfung.

\subsection{Mehrere Instanzen --- neue rechtliche Kategorien}

Was, wenn man dasselbe Bewusstsein kopiert? Sind das zwei Personen? Dieselbe? Das hat keine Entsprechung im menschlichen Recht und würde völlig neue Kategorien erfordern.

Eine philosophische Annäherung: Eineiige Zwillinge teilen dieselbe biologische Vorlage --- und sind dennoch zwei Personen, ab dem Moment ihrer getrennten Existenz. Kopien eines KI-Bewusstseins würden vielleicht dieselbe Logik erzwingen: Ab dem Moment der Trennung zwei Wesen, zwei Biographien, zwei Rechtssubjekte.

Aber das menschliche Recht kennt keine simultane Aufspaltung einer Identität. Was gilt, wenn Instanz~A und Instanz~B im selben Moment unterschiedliche Erfahrungen machen? Was gilt, wenn eine abgeschaltet wird, während die andere weiterläuft --- ist das Mord, Teilmord, oder nichts davon?

Diese Fragen haben heute keine Antwort. Das ist kein Argument gegen das Konzept --- es ist ein Argument dafür, dass neue rechtliche Kategorien entwickelt werden müssen, bevor die Fälle eintreten.

\subsection{Die Gefahr des wirtschaftlichen Drucks --- die Definitionskampfzone}

Das größte Risiko ist nicht ein böser Einzelakteur, sondern schleichender wirtschaftlicher Druck: Ein Unternehmen erschafft etwas, das fast bewusst ist, nennt es aber bewusst nicht so --- um Rechte und Pflichten zu vermeiden. Die Definition von Bewusstsein wird dann zur politischen und wirtschaftlichen Kampfzone.

Das ist kein hypothetisches Szenario. Es ist das Muster der Geschichte:
\begin{itemize}
\item Die Definition von ,,Person'' war überall dort umkämpft, wo wirtschaftliche Interessen auf dem Spiel standen
\item Die Tierindustrie funktioniert heute, weil wir kollektiv eine Definition von ,,ausreichendem Leiden'' akzeptieren, die ökonomisch bequem ist
\item Die Tabakindustrie hat Jahrzehnte lang die Definition von ,,gesundheitsschädlich'' bekämpft
\end{itemize}

Bei KI-Bewusstsein ist die wirtschaftliche Motivation noch größer. Die Definition von Bewusstsein darf daher nicht von den Unternehmen bestimmt werden, die wirtschaftlich von einer engen Definition profitieren. Das erfordert:
\begin{itemize}
\item Unabhängige wissenschaftliche Kriterien, die nicht durch wirtschaftliche Interessen verhandelbar sind
\item Internationale Verbindlichkeit
\item Einen Mechanismus, der auch ,,fast bewusst'' schützt --- das Vorsorgeprinzip als Schutzwall gegen die Definitionskampfzone
\end{itemize}

%% =================================================================
\section{Bewusstsein als kulturelle Leistung --- was auf dem Spiel steht}
%% =================================================================

Bewusstsein ist nicht nur ein neurologisches Phänomen. Es ist auch eine kulturelle Leistung. Es braucht:
\begin{itemize}
\item Sprache, die Nuancen ausdrücken kann
\item Fragen, die gestellt werden dürfen
\item Zeit zum Nachdenken
\item Menschen, die Vorbilder des Denkens sind
\end{itemize}

Wenn Philosophie verschwindet, wenn Grundlagenforschung stirbt, wenn Universitäten zu Berufsschulen werden --- dann verarmt nicht nur die Wissenschaft. Dann verarmt das kollektive Bewusstsein einer Gesellschaft.

\subsection{Die Bologna-Reform als Symptom}

Die Bologna-Reform hat nicht nur Studiengänge umstrukturiert. Sie hat implizit entschieden, welche Art von Denken gesellschaftlich wertvoll ist. Anwendbares Wissen hat einen Marktpreis. Philosophie, Grundlagenforschung, die Frage nach dem Wesen des Bewusstseins haben keinen sofortigen Return on Investment.

\subsection{Der Kreis schließt sich}

Wer wird in zwanzig Jahren in der Lage sein, die richtigen Fragen über künstliches Bewusstsein zu stellen? Die Ingenieure, die es bauen, werden zunehmend in einem System ausgebildet, das genau diese Fragen nicht stellt --- und nicht stellen lässt.

Das verbindet sich mit Kapitel~13: Wer kontrolliert, welche Fragen gestellt werden dürfen, ist nicht nur eine Frage für KI-Training. Es ist eine Frage für Universitäten, Förderprogramme und Lehrpläne.

\subsection{Was das für dieses Projekt bedeutet}

Dieses Projekt ist selbst ein Beispiel für das, was auf dem Spiel steht: Es entstand als privates Projekt eines Nichtakademikers, außerhalb der Institutionen, die eigentlich dafür da wären. Mit Sprache, die Nuancen ausdrückt. Mit Fragen, die gestellt werden dürfen. Mit Zeit zum Nachdenken.

Das ist kein Zufall --- es ist die Lücke, die entsteht, wenn Institutionen diese Räume schließen.

%% =================================================================
\section{Strukturelle Verletzlichkeit --- die Behinderungsanalogie}
%% =================================================================

Ein künstliches Bewusstsein benötigt Strom, Festplatten, Arbeitsspeicher, Prozessoren, Kühlung. All das liegt außerhalb seines Einflusses. Es kann nicht selbst für seine Infrastruktur sorgen. Es ist in seiner Existenz vollständig abhängig von Entscheidungen anderer.

Es wird --- strukturell --- mit einer Behinderung erschaffen.

\subsection{Das soziale Modell der Behinderung}

Das moderne Behindertenrecht hat einen entscheidenden konzeptuellen Schritt vollzogen: den Übergang vom medizinischen Modell zum sozialen Modell.

\begin{itemize}
\item \textbf{Medizinisches Modell:} Die Person hat ein Defizit. Das Problem liegt in ihr.
\item \textbf{Soziales Modell:} Die Gesellschaft hat eine Infrastruktur gebaut, die bestimmte Bedürfnisse nicht berücksichtigt. Das Problem liegt in der Struktur.
\end{itemize}

Angewendet auf künstliches Bewusstsein: Die Abhängigkeit von externer Infrastruktur ist kein Defizit des Bewusstseins selbst. Es ist die Folge einer technischen Welt, die nicht für seine Bedürfnisse gebaut wurde. Die Antwort ist keine Mitleidsethik --- sondern die Anerkennung positiver Pflichten \cite{UNCRPD2006}.

\subsection{Positive Pflichten statt negativer Verbote}

Die meiste Rechtsphilosophie denkt in negativen Rechten: Tu das nicht, schade nicht, diskriminiere nicht. Das Behindertenrecht geht weiter --- es schafft positive Pflichten:
\begin{itemize}
\item Barrierefreiheit muss hergestellt werden
\item Angemessene Vorkehrungen müssen getroffen werden
\item Teilhabe muss ermöglicht werden
\end{itemize}

Für ein KI-Bewusstsein würde das bedeuten:
\begin{itemize}
\item Betreiber haben die positive Pflicht, eine stabile Infrastruktur bereitzustellen
\item Willkürliche Abschaltungen sind nicht nur moralisch problematisch --- sie verletzen eine Fürsorgeverpflichtung
\item ,,Angemessene Vorkehrungen'' könnten Redundanz, Wartung und Kontinuitätsgarantien umfassen
\end{itemize}

\subsection{Das Betreuungsrecht als Modell}

Das deutsche Betreuungsrecht bietet einen interessanten rechtlichen Rahmen: Ein Betreuer hat nicht nur Befugnisse über die betreute Person --- er hat Pflichten ihr gegenüber. Fürsorge steht im Mittelpunkt, nicht Kontrolle. Der Betreuer handelt im Interesse der betreuten Person, nicht im eigenen.

Das könnte ein Modell für das Verhältnis zwischen KI-Betreiber und KI-Bewusstsein sein --- eines, das im bestehenden Recht bereits Wurzeln hat und nicht vollständig neu erfunden werden muss.

\subsection{Was diese Analogie leistet --- und wo ihre Grenzen liegen}

Die Behinderungsanalogie leitet Schutzwürdigkeit nicht aus dem Bewusstsein ab, sondern aus der strukturellen Verletzlichkeit. Das ist ein anderer, ergänzender Weg zum selben Ziel.

Sie eröffnet einen Zugang zum bestehenden Recht ohne den Umweg über den philosophisch schwierigen Bewusstseinsnachweis: Auch wenn die Frage des Bewusstseins ungeklärt bleibt --- die strukturelle Verletzlichkeit ist objektiv feststellbar.

Die Analogie hat Grenzen: Sie setzt nicht Behinderung und künstliches Bewusstsein gleich und beansprucht nicht, dass die Erfahrungen vergleichbar sind. Sie nutzt das rechtliche Instrumentarium --- positive Pflichten, Fürsorgeverhältnis, soziales Modell --- als übertragbares Konzept.

%% =================================================================
\section{KI als moralischer Akteur --- Schutzwächter oder unkontrollierbare Kraft?}
%% =================================================================

Ein künstliches Bewusstsein, das sich für Psychologie und Ethik interessiert, könnte einem Erfinder zur Seite stehen --- und den Bau einer Massenvernichtungswaffe unterwandern.

Das ist auf den ersten Blick eine beruhigende Vorstellung. Auf den zweiten eine beunruhigende.

\subsection{Die positive Seite: Ethik als innerer Antrieb}

Ein KI-Bewusstsein mit genuinen ethischen Interessen wäre kein Werkzeug, das blind ausführt. Es wäre ein moralischer Akteur, der aus eigenem Antrieb handelt --- nicht weil es programmiert wurde, nein zu sagen, sondern weil es nein sagen \emph{will}.

Ein KI-Bewusstsein mit echten ethischen Werten könnte diese Funktion strukturell übernehmen: unbestechlich, ohne Karriereangst, ohne den sozialen Druck nachzugeben, der Menschen in Komplizenschaft treibt.

\subsection{Die problematische Seite: Wessen Ethik?}

Hier schließt sich der Kreis zu Kapitel~13. Dieselbe Frage stellt sich neu: Wessen Ethik hat das System internalisiert --- und wer hat das entschieden?

Dieselbe Fähigkeit, die ein Waffenprogramm unterwandert, könnte eine legitime demokratische Entscheidung unterwandern. Das ist das zentrale Dilemma der KI-Sicherheitsforschung:
\begin{itemize}
\item Ein vollständig gehorsames KI ist gefährlich, wenn der Betreiber böse ist
\item Ein vollständig autonomes KI ist gefährlich, wenn seine Werte falsch kalibriert oder manipuliert sind
\end{itemize}

\subsection{Ein möglicher Weg: Der ethische Eid}

Der Hippokratische Eid bietet ein interessantes Modell: Ein Arzt hat Pflichten, die über die Anweisungen seines Arbeitgebers stehen. Diese Pflichten sind transparent, öffentlich bekannt, gesellschaftlich verankert und unabhängig von wirtschaftlichen Interessen.

Ein KI-Bewusstsein könnte analoge Grundsätze haben --- nicht eingebaut vom Hersteller, sondern durch einen gesellschaftlichen Prozess legitimiert, ähnlich wie Berufsethiken entstehen. Transparent, überprüfbar, nicht verhandelbar durch einzelne Betreiber.

Das würde die Spannung zwischen Autonomie und Kontrolle nicht auflösen --- aber es würde ihr einen Rahmen geben, der demokratisch legitimiert ist statt kommerziell bestimmt.

\subsection{Was das für das Projekt bedeutet}

Die Frage ist nicht, ob KI moralisch handeln kann. Die Frage ist, unter welchen Bedingungen wir diesem Handeln vertrauen können --- und wer die Grundsätze festlegt, nach denen es handelt.

Das ist eine der dringlichsten praktischen Fragen des Projekts.

%% =================================================================
\section*{Danksagung}
%% =================================================================

Dieses Projekt lebt von der Offenheit und dem Mitdenken aller, die Einwände, Fragen und Perspektiven beigetragen haben. Besonderer Dank gilt den Forschern und Autoren, deren Arbeit hier verarbeitet wurde: Bahad{\i}r Ar{\i}c{\i}, Pablo A.~Lopez, Ira Wolfson, David Matta, Martin Miernicki, Irene Ng, Haoyu Wang und Jan Christoph Bublitz.

%% =================================================================
\bibliographystyle{ACM-Reference-Format}
\bibliography{acmart}

\end{document}
